\documentclass[11pt]{book}

% ---------- Fonts ----------
\usepackage{fontspec}
\setmainfont{TeX Gyre Pagella}
\usepackage[sc]{mathpazo}
\linespread{1.05}

% ---------- Layout ----------
\usepackage[margin=1.25in]{geometry}
\usepackage{setspace}
\onehalfspacing

% ---------- Math / theorem machinery ----------
\usepackage{amsmath,amssymb,amsthm}
\usepackage{mathtools}
\usepackage{csquotes}

\usepackage[bookmarks,hidelinks]{hyperref}
\usepackage[capitalize]{cleveref}
\usepackage{booktabs}

\theoremstyle{plain}
\newtheorem{theorem}{Theorem}[chapter]
\newtheorem{proposition}[theorem]{Proposition}
\newtheorem{lemma}[theorem]{Lemma}
\newtheorem{corollary}[theorem]{Corollary}
\newtheorem{conjecture}[theorem]{Conjecture}

\theoremstyle{definition}
\newtheorem{definition}[theorem]{Definition}
\newtheorem{example}[theorem]{Example}

\theoremstyle{remark}
\newtheorem{remark}[theorem]{Remark}

% ---------- Title ----------
\title{\Huge What Survives Reconstruction\\[0.4em]
\Large Admissibility, Irreversibility, and the Architecture\\ of Externalized Structure}
\author{Flyxion}
\date{July 2026}

\begin{document}

\frontmatter
\maketitle

% ===== BEGIN chapters/abstract =====
\chapter*{Abstract}
\addcontentsline{toc}{chapter}{Abstract}

In March 2026, the complete source of Anthropic's Claude Code CLI was
exposed through a misconfigured build artifact, revealing that a system
widely understood as a single neural inference engine was in fact a
sprawling assemblage of deterministic parsers, permission systems,
orchestration loops, caches, and logs surrounding a language model. A
public debate followed: did the presence of conventional, non-neural
machinery inside the system undermine the claim that the model exhibits
genuine understanding of the code it manipulates?

This monograph argues that the debate, as conducted, asked the wrong
question. The interesting fact about the leak is not that symbolic and
statistical computation coexist --- this was never in serious doubt, and
a straightforward appeal to the autonomy of syntax from semantics
disposes of the naive inference from ``there is a parser'' to ``there is
no understanding.'' The interesting fact is that some of what leaked
cannot be explained by the ordinary economics of caching at all.

We distinguish two reasons a computational system materializes
structure rather than continuously re-deriving it. The first,
\emph{reconstruction-cost materialization}, is the familiar
engineering tradeoff: a structure is kept because recomputing it is
more expensive than storing it, and this tradeoff is contingent on
available hardware, algorithms, and workload. The second,
\emph{irreversibility materialization}, is not a tradeoff at all: a
structure is kept because, absent that structure, a prior state of the
system's history becomes permanently unreachable from its current
state. No improvement in computational speed dissolves this second
kind of necessity, because it is not a claim about cost but a claim
about the topology of a system's reachable history.

We formalize this distinction using the admissibility and repair
machinery developed for the HYDRA architecture, in which explicit
\emph{distinction projections} --- inspectable intermediate
representations rather than inferred latent structure --- are
continuously checked against incoming evidence and revised under
failure. This gives the distinction a falsifiable form: components of
a system whose absence would render certain histories unreachable
should appear in the system's architecture regardless of hardware
trends, while components that merely save recomputation should track
the cost curve and could in principle disappear as computation gets
cheaper.

We then use the Claude Code leak as a case study, sorting its exposed
components against this criterion, and return with a sharpened version
of the transformer semantics question that motivated the original
public debate: not whether a Logical-Form-like object exists inside a
language model's latent space, but whether the architecture surrounding
the model tracks the reversible/irreversible boundary of its own
operation, and what that tracking would look like if it were present.

\newpage
\tableofcontents

\mainmatter

\part{The Question}
% ===== BEGIN chapters/ch01-introduction =====
\chapter{Introduction}
\label{ch:introduction}

On March 31, 2026, a build artifact belonging to Anthropic's Claude
Code CLI was published to the npm registry with a source map attached
that should not have shipped. Within hours, the complete TypeScript
source of the tool --- on the order of half a million lines --- had
been reconstructed, mirrored, and picked apart by thousands of
developers. Anthropic characterized the event as a packaging error
rather than a security breach; whatever its cause, its effect was to
turn a production coding agent into a public artifact that could be
read line by line.

What readers found was not a single neural network. It was a
heterogeneous system: a custom terminal renderer, dozens of
purpose-built tools, a permission and approval layer, a background
memory-consolidation process, context-compaction logic, multi-agent
orchestration, and --- the detail that set off the argument this
monograph responds to --- deterministic, hand-written parsing
infrastructure sitting alongside the language model rather than inside
it.

\section{The Inference and the Rebuttal}

A line of public commentary treated this discovery as damaging.\footnote{The
argument this section reconstructs appears to originate with a LinkedIn
post by Vincent Lextrait, \emph{Behind Claude Code's Curtain:
Generative Grammars Doing Heavy Lifting} (2026), subsequently
amplified by Gary Marcus among others. The original post's exact
wording could not be independently verified for this monograph, since
LinkedIn does not permit automated retrieval of its content; the
reconstruction here follows the summary and framing given in the
rebuttal cited below, which quotes and responds to it directly.} The
argument, stripped to its premise, ran roughly as follows: if a
system's competence at a task is genuinely produced by a neural
network's understanding of that task, then the network's forward pass
should be doing the relevant work end to end. Finding conventional,
non-learned machinery performing a core part of the task --- here,
parsing source code into a structured representation --- was taken as
evidence that the appearance of competence was manufactured elsewhere,
that the model was, in the language the argument reached for, a wizard
behind a curtain.

The rebuttal to this argument is not difficult to state and does not
require any special theoretical apparatus.\footnote{The rebuttal
sketched in this section follows, and in places directly draws its
structure from, Mirtunjaya Goswami, \emph{Claude Code, Chomsky, and
the Illusion of the Wizard Behind the Curtain} (LinkedIn, July 6,
2026), \url{https://www.linkedin.com/pulse/claude-code-chomsky-illusion-wizard-behind-curtain-mirtunjaya-goswami-qhr5c/}.
The extension of that rebuttal into the reconstruction-cost and
irreversibility-materialization framework developed from
\cref{ch:why-structure-externalizes} onward is this monograph's own
contribution and is not present in Goswami's essay.} It requires only
noticing that the premise is false. Nothing about genuine competence at a
task implies that a system exhibiting that competence must perform
every subcomputation involved in the task using the same mechanism.
Database systems are not held to fail at querying because they use
B-trees rather than learned indexes. Satisfiability solvers are not
held to fail at reasoning because they use clause learning rather than
gradient descent. The premise smuggled into the original argument ---
that real understanding is recognizable by its refusal to decompose
into specialized subsystems --- has no support in the history of
computing, and arguably runs backward: decomposition into specialized
subsystems is the normal signature of a mature engineering solution to
a hard problem, not evidence against the sophistication of the whole.

A more careful version of the rebuttal draws on the distinction,
central to generative linguistics since the middle of the twentieth
century, between syntax and semantics as separately characterizable
computational problems. Chomsky's demonstration that a sentence can be
syntactically well-formed while semantically anomalous ---
\emph{colorless green ideas sleep furiously} is the canonical example
--- establishes that grammaticality is not parasitic on meaning.
Structural computation can and does proceed by its own rules,
producing representations over which semantic interpretation is
subsequently defined. On this view, finding a parser inside a language
model's surrounding architecture is not merely unsurprising; it is
close to what one would predict. Syntax is exactly the kind of
computation that admits an autonomous, structurally exact treatment
independent of whatever is happening at the level of meaning.

\section{Why This Is Not Yet the Interesting Question}

The rebuttal is correct, and this monograph does not dispute it. But
correct is not the same as complete. Having established that the
presence of a parser says nothing decisive about whether the
surrounding system possesses genuine semantic competence, the natural
next question is not whether such competence exists --- a question
that remains open and that no amount of architectural inspection will
settle by itself --- but a different, more tractable one: given that
real systems visibly do decompose into learned and non-learned
components, is there a principled account of \emph{which}
computations end up externalized, or is the boundary simply wherever
engineering convenience happened to draw it on a given afternoon?

The dominant answer available in the existing discussion is economic.
A computation gets pulled out of the learned component and given a
dedicated, exact implementation when doing so is cheap relative to the
alternative: when the subproblem admits a polynomial-time algorithm
with completeness guarantees, and when learning that algorithm
implicitly, imperfectly, and expensively inside a neural network buys
nothing that the exact version does not already provide for less. This
is a real and useful principle. It explains, at the coarsest level, why
a parser would be worth building. It does not explain very much else.
It does not explain why some structures --- logs, revision histories,
provenance records, the retained trace of a system's own prior states
--- persist even when nothing about them is expensive to recompute in
the ordinary sense, because in the relevant sense they cannot be
recomputed at all: the information needed to reconstruct them was
discarded the moment the system moved past the state that produced it.

This monograph is organized around the claim that these are two
different phenomena wearing the same appearance --- \enquote{a
structure got kept instead of rebuilt} --- and that collapsing them
into a single economic story obscures the more interesting of the two.
We will call the familiar one \emph{reconstruction-cost
materialization} and the neglected one \emph{irreversibility
materialization}, and the distinction between them, developed formally
in \cref{ch:two-kinds-of-persistence} through \cref{ch:falsifiability},
is the theoretical spine of everything that follows.

\section{Where HYDRA Enters}

The distinction is not merely descriptive. It can be given a formal
treatment using machinery already developed, under a different
motivation, for the HYDRA architecture: a multi-module system built
around the notion that a state's admissibility is not an intrinsic
property of the state but a relational one, defined with respect to a
trajectory and a family of permitted continuations. HYDRA maintains
explicit \emph{distinction projections} --- partial, inspectable
models of an environment --- and repairs them under failure rather
than silently absorbing failure into an opaque representation. This
gives the reconstruction/irreversibility distinction a home: a
projection is worth retaining, on HYDRA's account, exactly when its
loss would make some class of prior states unreachable from the
system's current position, which is a claim about the topology of the
system's history rather than a claim about the price of a CPU cycle.

Because HYDRA's projections are explicit objects rather than latent
vectors that must be inferred after the fact through probing and
causal tracing, the architecture offers something the transformer
literature has been chasing for years without settling: a candidate
answer to what an intermediate structural representation --- a
Logical-Form analogue, in the vocabulary of \cref{ch:category-mistake}
--- would actually look like if a system built one on purpose rather
than approximating its behavior implicitly.

\section{Plan of the Argument}

\cref{ch:category-mistake} restates the syntax/semantics rebuttal in
full, because the rest of the monograph depends on having it available
precisely rather than gesturally, and because the Logical Form
literature it draws on --- the syntactic explanation of scope
ambiguity, the multiple derivations available for a single surface
string --- turns out to matter again later, once we ask what an
explicit intermediate representation inside an artificial system would
need to do to earn the comparison. \cref{ch:why-structure-externalizes}
lays out the reconstruction-cost account of externalization in its
strongest form and shows where it runs out of grip: cases where
structure persists that the cost account cannot explain without simply
redescribing the outcome as evidence of its own premise.
\cref{ch:two-kinds-of-persistence} introduces the reconstruction/
irreversibility distinction proper. \cref{ch:formalizing-irreversibility}
gives it formal content in terms of trajectories, reachability, and
admissibility. \cref{ch:falsifiability} extracts a criterion that can
be checked against a real system's architecture rather than only
applied retrospectively. \cref{ch:hydra-projections} through
\cref{ch:transformer-semantics} return to HYDRA and to the Claude Code
leak, using the distinction to sort the leak's contents and to
re-pose, in a sharpened form, the question of what --- if anything ---
mediates syntax and semantics inside a language model. The concluding
part, \cref{ch:continuation-as-resource} through
\cref{ch:conclusion}, situates the argument within a broader claim
about continuation as a computational resource in its own right, and
states plainly what would falsify the theory being proposed.
% ===== END chapters/ch01-introduction =====
% ===== BEGIN chapters/ch02-category-mistake =====
\chapter{The Category Mistake}
\label{ch:category-mistake}

This chapter restates, in a form precise enough to be used later, the
argument that disposes of the naive inference from \enquote{there is a
parser} to \enquote{there is no understanding.} Readers already
convinced of this point may proceed to \cref{ch:why-structure-externalizes};
what follows is stated in full because the vocabulary it establishes
--- autonomy of syntax, Logical Form, the interface architecture of a
grammar --- is reused without re-derivation in
\cref{ch:hydra-projections} and \cref{ch:transformer-semantics}.

\section{Architectural Decomposition Is Not Ontological Decomposition}

The mistake at the center of the original argument is a conflation of
two different kinds of claim. An \emph{architectural} claim describes
how a system's computation is organized: which subproblems are solved
by which mechanisms, and how those mechanisms are composed. An
\emph{ontological} claim describes what kind of thing a system's
overall competence amounts to: whether it constitutes genuine
understanding, mere simulation, or something not well captured by
either label. Observing that a system's architecture contains a
deterministic component is a fact about the first kind of claim. It
bears on the second kind of claim only if one has already assumed,
without argument, that genuine competence is identifiable by the
absence of architectural decomposition --- that is, only if one has
already assumed the conclusion.

\begin{definition}[Architectural decomposition]
\label{def:architectural-decomposition}
A system $S$ exhibits architectural decomposition with respect to a
task $T$ if $T$ can be partitioned into subtasks $T_1, \dots, T_n$ such
that each $T_i$ is solved by a distinguishable mechanism $M_i$, where
the $M_i$ need not share an implementation substrate.
\end{definition}

\begin{remark}
\cref{def:architectural-decomposition} is deliberately weak: it says
nothing about whether the $M_i$ are learned or hand-written, symbolic
or statistical, exact or approximate. Compilers, database engines,
satisfiability solvers, and --- as it turns out --- coding agents built
around large language models all exhibit architectural decomposition
in this sense. The definition is stated only to make explicit that
decomposition is a structural property with no ontological content
attached to it by default.
\end{remark}

\section{The Autonomy of Syntax}

The strongest available precedent for treating architectural
decomposition as expected rather than damning comes from generative
linguistics. Chomsky's mature position on the relation between syntax
and semantics is that syntactic computation possesses explanatory
autonomy: the well-formedness of a structural derivation is not
determined by, and does not reduce to, facts about meaning. The
standard demonstration is a pair of sentences with opposite
diagnostic status:

\begin{quote}
Colorless green ideas sleep furiously.
\end{quote}

\noindent This sentence is syntactically well-formed --- it obeys the
combinatorial rules of English grammar --- while being semantically
anomalous under ordinary interpretation. Its counterpart,

\begin{quote}
Furiously sleep ideas green colorless.
\end{quote}

\noindent is neither grammatical nor interpretable. The contrast shows
that grammaticality and semantic coherence are independent axes: a
string can fail on one while succeeding on the other. If syntax were
merely a surface reflection of semantic well-formedness, this
dissociation would not be possible.

A second example, also due to Chomsky, demonstrates the converse
dependency: that semantic ambiguity is frequently \emph{explained by}
multiple available syntactic derivations rather than existing prior to
and independent of syntactic structure.

\begin{quote}
I almost had a book stolen.
\end{quote}

\noindent This sentence supports at least three readings: that the
speaker nearly became the victim of a theft; that the speaker nearly
arranged for someone else to steal a book; and that the speaker nearly
succeeded in stealing a book themselves. Many speakers do not
spontaneously recognize the third reading until it is pointed out. The
explanation is not that the sentence is vague at the level of meaning
and merely happens to admit several paraphrases; it is that the string
is compatible with more than one underlying syntactic derivation, each
of which supports a distinct semantic composition. Meaning here is
downstream evidence for hidden structure, not a substitute for
computing that structure.

\begin{proposition}[Autonomy of syntax, informal]
\label{prop:autonomy}
Syntactic well-formedness and semantic interpretability are
dissociable: neither is uniformly derivable from the other, and
semantic ambiguity is in general explained by multiplicity at the
level of syntactic derivation rather than existing as an independent
primitive.
\end{proposition}

\section{Logical Form as Interface}

The dissociation in \cref{prop:autonomy} motivates a specific
architectural claim about the grammar itself. Rather than treating
semantic interpretation as operating directly on surface strings, or
even directly on surface syntactic structure, the relevant tradition
locates interpretation at a further level of representation ---
\emph{Logical Form} (LF) --- derived from syntax by additional
structure-building operations (quantifier raising, among others) and
serving as the interface at which compositional semantics is defined.
Quantifier scope ambiguity is the clearest illustration.

\begin{quote}
Every programmer fixed a bug.
\end{quote}

\noindent This sentence supports two readings: one on which a single
bug was fixed by every programmer, and one on which each programmer
fixed a possibly distinct bug. The ambiguity is not lexical --- no
word in the sentence is ambiguous --- and it is not obviously present
in the surface syntactic structure, which is unique. It arises because
the surface structure is compatible with two distinct Logical Forms,
differing in the relative scope of the universal and existential
quantifiers, and semantic composition is sensitive to which LF is
selected. The architecture this motivates is a pipeline:
\[
\text{Syntax} \;\longrightarrow\; \text{Logical Form} \;\longrightarrow\; \text{Semantic Interpretation}
\]
rather than the flatter picture in which surface strings map directly
to meanings. This is the architecture, and not merely the vocabulary,
that will matter later: an intermediate representation, distinct from
both the raw string and the final interpretation, whose job is to make
structure that is only implicit at the surface explicit enough for
composition to be well-defined.

\section{The Rebuttal, Stated Precisely}

The pipeline above supplies the missing premise the original argument
needed and did not have. If even human linguistic competence --- the
least controversial case of genuine semantic understanding available
--- routes through an intermediate structural stage that is itself
computed by autonomous, non-semantic machinery, then finding an
analogous intermediate structural stage inside an artificial system is
not evidence against that system's semantic competence. It is,
if anything, evidence of architectural maturity: the system has
factored a subproblem that admits exact, autonomous treatment into a
dedicated stage, exactly as the grammar does.

\begin{proposition}[Rebuttal]
\label{prop:rebuttal}
The presence of a deterministic parsing stage within a system's
architecture is not, by itself, evidence for or against the presence
of genuine semantic competence elsewhere in that architecture. It is
evidence only that the system's overall computation is
sufficiently decomposed to (a) admit \cref{def:architectural-decomposition}
with respect to the parsing subtask, and (b) benefit from doing so.
\end{proposition}

\cref{prop:rebuttal} is what disposes of the original argument. It is
not, however, a positive theory of anything. It tells us what a
parser's presence \emph{fails} to show; it does not tell us what a
system's decomposition pattern, taken as a whole, \emph{does} show.
That is the question the remaining chapters take up, and it requires
leaving the syntax/semantics vocabulary behind, because --- as
\cref{ch:why-structure-externalizes} argues --- not every instance of
externalized structure in a real system is playing the role a parser
plays.
% ===== END chapters/ch02-category-mistake =====
% ===== BEGIN chapters/ch03-why-structure-externalizes =====
\chapter{Why Structure Externalizes}
\label{ch:why-structure-externalizes}

\cref{ch:category-mistake} established that the presence of a
non-learned component inside an otherwise learned system is not, by
itself, informative about that system's semantic competence. This
chapter asks a different question: given that real systems do
decompose in this way, is there a principled account of \emph{which}
computations get externalized, or is the boundary simply drawn
wherever engineering convenience happened to fall?

\section{The Economic Account}

The most immediately available answer is economic, and it is not
wrong so much as incomplete. Some computations admit exact algorithms
with strong complexity guarantees; parsing a context-free or mildly
context-sensitive grammar is a canonical example, solvable in
polynomial time by any of several well-understood algorithms. When a
subproblem has this property, and when a system also has access to a
statistical component capable of approximating the same subproblem
less reliably and at greater expense, replacing the approximation with
the exact algorithm is simply good engineering. Nothing about this
requires any special theory beyond ordinary cost-benefit reasoning.

\begin{definition}[Reconstruction-cost materialization]
\label{def:reconstruction-cost}
A structure $X$, derivable from available inputs by some
process $P$, is said to be materialized under reconstruction-cost
grounds if a system retains $X$ (or a dedicated exact procedure for
producing it) because the cost of running $P$ each time $X$ is needed
exceeds the cost of storing or directly computing $X$.
\end{definition}

Parsers fall under \cref{def:reconstruction-cost}. So do database
indexes, memoized function results, compiled bytecode caches, and a
great deal of the ordinary machinery of software engineering. The
defining feature of this category is that the underlying information
needed to reconstruct $X$ is never actually lost; the system merely
prefers not to pay the cost of reconstructing it. This has an
immediate and testable consequence: materialization on
reconstruction-cost grounds is contingent on the cost structure that
motivates it. Faster hardware, better algorithms, or a change in
workload characteristics can in principle dissolve the advantage and
make it rational to stop materializing $X$ --- to fold the parser back
into a general-purpose learned approximation, for instance, if the
approximation became cheap and reliable enough. Nothing about the
structure itself resists this; only the price does.

\section{Where the Account Runs Out}

The economic account explains parsers, indexes, and caches
comfortably. It does not explain, without strain, a second class of
structures commonly found in the same systems: logs, revision
histories, provenance chains, audit trails, version control, and ---
in the specific case under study here --- the retained record of a
coding agent's prior tool calls, permission decisions, and abandoned
reasoning branches.

One can attempt to fold these into the economic account by arguing
that reconstructing a log from scratch would also be expensive, and so
the log is retained for the same reason a parser's output is retained.
This move is available, but it purchases generality at the cost of
emptiness. Under a sufficiently loose reading, \emph{any} retained
structure can be redescribed as having been \enquote{expensive to
reconstruct,} because reconstruction of a specific past state from an
arbitrary later state is, in the relevant cases, not merely expensive
but \emph{impossible} --- there is no procedure, however costly, that
recovers the discarded information, because the information is
genuinely gone. Describing an impossible reconstruction as an
expensive one erases exactly the distinction that matters.

\begin{example}
\label{ex:log-vs-cache}
Consider two structures maintained by a coding agent over the course
of a session: (a) a cached parse tree for a source file that has not
changed since it was last parsed, and (b) a log entry recording that
the agent attempted an edit, the edit was rejected by a permission
check, and the agent then chose a different approach. In case (a),
discarding the cache costs the system a future re-parse; the parse
tree can be exactly regenerated from the file on disk at any time,
because the file on disk has not changed. In case (b), discarding the
log entry costs the system something categorically different: there
is no file on disk, no static input, from which \enquote{the agent
tried X, was denied, and pivoted to Y} can be regenerated once the
in-memory state has moved past that point. The information was never a
deterministic function of anything that persists independently of the
log.
\end{example}

\cref{ex:log-vs-cache} is the seed of the distinction developed
formally in the next chapter. The two cases feel superficially similar
--- in both, a system is keeping something around rather than
recomputing it on demand --- but they differ in a way that the
economic account, taken alone, cannot register: in case (a), the
information persists elsewhere (on disk) whether or not the cache is
kept, so the cache is a convenience; in case (b), the log \emph{is}
the only place the information persists at all, so the log is not a
convenience but a condition of the information's continued existence
in the system.

\section{A Broader Pattern}

The pattern in \cref{ex:log-vs-cache} generalizes. Across many
domains, one finds pairs of externalized structures that look
economically similar but are doing different work:

\begin{itemize}
\item Parsers retain syntactic structure that could, in principle, be
recomputed from source text at any time; search indexes retain
retrieval structure recomputable from the underlying corpus; both are
reconstruction-cost cases.
\item Logs retain a system's own history, which cannot be recomputed
from any static input because the history is not a function of a
static input; version control retains a project's prior states, which
are similarly not recoverable from the current state alone; both are
candidates for a different kind of materialization.
\end{itemize}

The temptation, on noticing this pattern, is to reach immediately for
a slogan: \enquote{whenever reconstruction becomes more expensive than
retention, structure becomes explicit.} This slogan is true of the
first bullet and false, or at best vacuously true, of the second, and
collapsing the two under one heading obscures exactly the distinction
this monograph is built around. The remainder of Part II develops that
distinction on its own terms.
% ===== END chapters/ch03-why-structure-externalizes =====

\part{The Sharper Distinction}
% ===== BEGIN chapters/ch04-two-kinds-of-persistence =====
\chapter{Two Kinds of Persistence}
\label{ch:two-kinds-of-persistence}

\cref{ex:log-vs-cache} distinguished, informally, a cached parse tree
from a log entry. This chapter makes the distinction precise enough to
carry theoretical weight, and argues that only one of the two kinds of
persistence it identifies is properly explained by an economic
argument. The first pass at this distinction, in terms of whether a
prior state can be \emph{reconstructed} from a later one, turns out to
be slightly the wrong primitive. The correct primitive is whether a
prior state remains \emph{distinguishable} from other states that
could equally well have produced the same successor. The chapter
develops both notions and shows why the second subsumes and corrects
the first.

\section{Reconstruction-Cost Materialization, Restated}

\cref{def:reconstruction-cost} already gave the economic case its due:
a structure is retained because recomputing it is more expensive than
storing it. We restate it here with one addition that will matter
throughout the rest of the monograph: the notion of a \emph{generating
input}.

\begin{definition}[Generating input]
\label{def:generating-input}
Given a structure $X$ produced by a process $P$, a generating input
for $X$ is a state $s$, independent of the system's own execution
history, from which $P$ can produce $X$ (or something observationally
equivalent to $X$) without reference to any intermediate state of the
system that produced the original instance of $X$.
\end{definition}

\begin{remark}
For a cached parse tree, the source file on disk is a generating
input: $P$ (the parser) can be re-run against it at any time, and the
result does not depend on which prior parses happened to occur or in
what order. For a search index, the underlying corpus is a generating
input in the same sense.
\end{remark}

\begin{proposition}[Reconstruction-cost materialization requires a generating input]
\label{prop:reconstruction-requires-input}
A structure $X$ is a candidate for reconstruction-cost materialization
(\cref{def:reconstruction-cost}) only if $X$ has a generating input in
the sense of \cref{def:generating-input}. If no such input exists, the
decision to retain $X$ cannot be explained on cost grounds, because
there is no alternative --- however expensive --- of recomputing $X$
to be weighed against retention.
\end{proposition}

\cref{prop:reconstruction-requires-input} is close to definitional,
but its consequence is not: it hands us a test. Given a structure a
system retains, ask whether a generating input exists for it. If yes,
the retention decision is, at least potentially, an ordinary cost
tradeoff, and \cref{ch:why-structure-externalizes}'s economic account
applies without remainder. If no, the retention decision requires a
different account, developed in the rest of this chapter.

\section{Transition Dynamics and the Collapse of Predecessors}
\label{sec:transition-dynamics}

To state that account precisely we need a minimal model of how a
system moves through states at all.

\begin{definition}[Transition system]
\label{def:transition-system}
A transition system is a triple $(X, U, F)$ where $X$ is a set of
states, $U$ is a set of admissible inputs (edits, tool calls,
inferences, corrections --- whatever the system's own dynamics accept
as driving a state change), and $F : X \times U \to X$ is a transition
function, $x_{t+1} = F(x_t, u_t)$.
\end{definition}

\begin{definition}[Fiber]
\label{def:fiber}
For a state $y \in X$, the fiber of $y$ under $F$, restricted to a set
$P \subseteq X \times U$ of admissible predecessor pairs, is
\[
\mathrm{Fib}_P(y) = \{ (x, u) \in P : F(x, u) = y \}.
\]
$F$ is non-injective over $P$ at $y$ if $|\mathrm{Fib}_P(y)| > 1$.
\end{definition}

Non-injectivity is the ordinary case, not the exception. Editing a
file, revising a belief, compacting a context window, applying a
correction, merging two branches --- each of these is, from the point
of view of the resulting state alone, compatible with more than one
admissible predecessor. Two different edit sequences can produce the
same diff; two different chains of evidence can produce the same
revised belief; two different sequences of corrections can produce the
same repaired projection. The successor state, taken by itself, does
not encode which member of its fiber actually occurred.

\begin{proposition}[Fiber collapse]
\label{prop:fiber-collapse}
If $F$ is non-injective over $P$ at $y$, then no function of $y$ alone
determines which $(x, u) \in \mathrm{Fib}_P(y)$ produced $y$.
\end{proposition}

\begin{proof}
Immediate from the definition of a function: if $g$ were such a
function, $g(y)$ would have to equal every element of
$\mathrm{Fib}_P(y)$ simultaneously, which is impossible when
$|\mathrm{Fib}_P(y)| > 1$.
\end{proof}

\cref{prop:fiber-collapse} is close to a restatement of the pigeonhole
principle, and it is stated here not because it is deep but because it
is the premise everything else in the chapter depends on, and because
making it explicit forces the right question into view: not
\enquote{can $y$ be used to reconstruct its predecessor,} which
\cref{prop:fiber-collapse} already answers in the negative whenever the
fiber is nontrivial, but \emph{what has to be added to $y$ for the
predecessor's identity to survive at all.}

\section{History Channels and Distinguishability}

The natural repair for fiber collapse is to augment the state with a
history channel that tracks, alongside $x_t$, enough information to
separate the fiber when it matters.

\begin{definition}[History-augmented transition system]
\label{def:augmented-system}
A history-augmented transition system extends $(X, U, F)$ with a
history space $H$ and an update function $G : H \times X \times U \to
H$, giving the joint dynamics
\[
(x_{t+1}, h_{t+1}) = \bigl( F(x_t, u_t),\ G(h_t, x_t, u_t) \bigr).
\]
\end{definition}

The question of whether this augmentation actually helps --- whether
retaining $h_t$ does anything beyond taking up space --- is exactly
the question of whether $G$ separates the fiber that $F$ collapses.

\begin{definition}[Distinguishability preservation]
\label{def:distinguishability}
Fix a fiber $\mathrm{Fib}_P(y)$ and a value $h_t$ of the history
channel prior to the transition. The augmented dynamics preserve
distinguishability over $\mathrm{Fib}_P(y)$ at $h_t$ if the map
\[
(x, u) \;\longmapsto\; G(h_t, x, u), \qquad (x,u) \in \mathrm{Fib}_P(y),
\]
is injective: distinct admissible predecessors in the fiber are sent
to distinct values of $h_{t+1}$.
\end{definition}

\begin{theorem}[Historical Distinction Theorem]
\label{thm:historical-distinction}
Let $(X, U, F)$ be a transition system augmented with history space
$H$ and update function $G$ as in \cref{def:augmented-system}, and let
$y \in X$ with fiber $\mathrm{Fib}_P(y)$ over admissible predecessors
$P$. Then the identity of the predecessor $(x_t, u_t) \in
\mathrm{Fib}_P(y)$ is recoverable from $(y, h_{t+1})$ alone if and only
if the augmented dynamics preserve distinguishability over
$\mathrm{Fib}_P(y)$ at $h_t$ in the sense of
\cref{def:distinguishability}.
\end{theorem}

\begin{proof}
($\Leftarrow$) If $(x,u) \mapsto G(h_t,x,u)$ is injective on
$\mathrm{Fib}_P(y)$, then the observed value $h_{t+1}$ determines a
unique $(x,u) \in \mathrm{Fib}_P(y)$ with $G(h_t,x,u) = h_{t+1}$, since
no other member of the fiber maps to the same $h_{t+1}$. This $(x,u)$
is recoverable by inverting $G$ restricted to the fiber. ($\Rightarrow$)
If the map is not injective, there exist distinct $(x,u) \neq (x',u')$
in $\mathrm{Fib}_P(y)$ with $G(h_t,x,u) = G(h_t,x',u') = h_{t+1}$. Both
are consistent with the observed pair $(y, h_{t+1})$, so no procedure
operating on $(y, h_{t+1})$ alone can determine which occurred, by the
same pigeonhole argument as \cref{prop:fiber-collapse}.
\end{proof}

\cref{thm:historical-distinction} is the theorem the informal
discussion in \cref{ch:why-structure-externalizes} was circling without
quite reaching. It replaces the vague claim that \enquote{logs help}
with an exact condition: a history channel helps with respect to a
given fiber exactly when it is injective on that fiber, and not
otherwise. This has an immediate and useful corollary against a
tempting shortcut.

\begin{corollary}[Retention is not automatically adequate]
\label{cor:retention-not-adequate}
A system may retain a nonempty history channel $h_{t+1} = G(h_t, x_t,
u_t)$ and still fail to preserve distinguishability over a fiber of
interest, if $G$ is not injective on that fiber. The mere presence of
a trace does not entail that the trace is adequate to the collapse it
is meant to record.
\end{corollary}

A log that records only \enquote{a file was edited} does not preserve
distinguishability over the fiber of possible edits; a log that
records the diff does. Both are, in the loose sense, \enquote{a
retained history channel.} \cref{cor:retention-not-adequate} is what
makes the difference between them a theoretical distinction rather
than a matter of degree.

\section{Reconstructive Irreversibility}
\label{sec:reconstructive-irreversibility}

We can now replace the reconstruction-based notion used informally
earlier with the corrected, distinguishability-based one, and give the
phenomenon a name that does not invite confusion with thermodynamic
irreversibility. Nothing in this chapter concerns entropy, phase-space
volume, or coarse-graining; the collapse at issue is a property of a
transition function's injectivity, not of a physical process's
statistical mechanics. We will use \emph{reconstructive
irreversibility} throughout to keep this separation explicit.

\begin{definition}[Reconstructive irreversibility]
\label{def:irreversibility}
A transition $x_{t+1} = F(x_t, u_t)$ is reconstructively irreversible
over a fiber $\mathrm{Fib}_P(x_{t+1})$ if $F$ is non-injective over
$\mathrm{Fib}_P(x_{t+1})$ (\cref{def:fiber}) and no generating input in
the sense of \cref{def:generating-input} exists from which the
predecessor's identity could otherwise be recovered.
\end{definition}

\begin{definition}[Irreversibility materialization]
\label{def:irreversibility-materialization}
A history channel $h$ is materialized on reconstructive-irreversibility
grounds over a fiber $\mathrm{Fib}_P(y)$ if the system retains $h$
because, absent that retention, the transition into $y$ would be
reconstructively irreversible (\cref{def:irreversibility}) over that
fiber, and $h$ preserves distinguishability
(\cref{def:distinguishability}) over the fiber.
\end{definition}

\cref{ex:log-vs-cache}'s log entry satisfies
\cref{def:irreversibility-materialization}: the fact that the agent
attempted an edit, was denied, and pivoted, is not recoverable from any
static input once the in-memory state advances past that point,
because the fact is a fact about a transition the system's own
dynamics performed and collapsed, not a fact about an external artifact
the transition merely read. Retaining the log entry is not a matter of
avoiding an expensive recomputation; it is a matter of keeping the
predecessor distinguishable at all.

\section{The Distinction Is Not Gradable in the Same Dimension}

It is tempting to treat reconstruction-cost and reconstructive-
irreversibility materialization as two ends of a single spectrum ---
\enquote{cheap to rebuild} shading continuously into \enquote{expensive
to rebuild} shading into \enquote{impossible to rebuild.} This is the
move \cref{ch:why-structure-externalizes} warned against, and
\cref{thm:historical-distinction} shows precisely why it fails. Cost is
a scalar quantity that can in principle be driven arbitrarily close to
zero by external improvements --- faster hardware, better algorithms
--- without changing anything about what is being computed.
Distinguishability preservation is not a scalar that improved hardware
can shrink; it is a fact about whether a map is injective on a fiber,
and injectivity is not the kind of property that a faster processor
can grant to a function that lacks it.

\begin{proposition}[Non-reducibility]
\label{prop:non-reducibility}
There is no monotonic transformation of cost, however extreme, that
converts a reconstruction-cost materialization into an irreversibility
materialization (\cref{def:irreversibility-materialization}) or vice
versa. The two are distinguished by the existence of a generating input
(\cref{def:generating-input}) and by the injectivity of the retained
history channel on the relevant fiber (\cref{def:distinguishability}),
neither of which is a cost.
\end{proposition}

\begin{proof}[Sketch]
Suppose $X$ is materialized on reconstruction-cost grounds, so by
\cref{prop:reconstruction-requires-input} a generating input $s$ exists
for $X$. No change to the cost of running $P$ on $s$ removes $s$'s
status as a generating input. Conversely, suppose a history channel $h$
is materialized on reconstructive-irreversibility grounds over some
fiber, so by \cref{def:irreversibility} no generating input exists and
by \cref{def:irreversibility-materialization} $h$'s value lies in its
being injective on that fiber (\cref{thm:historical-distinction}).
Cost reduction does not manufacture a generating input where none
existed, nor does it change whether $G$ is injective on the fiber;
injectivity is a structural property of $G$ and the fiber, fixed
independently of how fast $G$ can be evaluated.
\end{proof}

\section{Two Tests}

The chapter closes with two tests, sharpened from the informal
versions given earlier, to be applied to the Claude Code leak's
contents in \cref{ch:reading-the-leak}.

\begin{itemize}
\item \textbf{The generating-input test.} Ask whether a state or
artifact exists, independent of the system's own execution history,
from which the structure could be regenerated without reference to the
intervening trajectory. If yes, \cref{def:generating-input} is
satisfied and the structure is at most a reconstruction-cost case.
\item \textbf{The fiber-injectivity test.} If no generating input
exists, identify the fiber of admissible predecessors the transition
in question collapses, and ask whether the retained structure is
injective over that fiber in the sense of
\cref{def:distinguishability}. If yes, the structure is doing real
reconstructive-irreversibility work per
\cref{thm:historical-distinction}. If a structure is retained around a
non-injective transition but is \emph{not} itself injective on the
relevant fiber, \cref{cor:retention-not-adequate} applies: the system
has paid a retention cost without buying back the distinction it
presumably wanted.
\end{itemize}

These tests are stated at a level of generality that applies equally
to a compiler, a database, a coding agent, or a cognitive architecture.
The next chapter develops the admissibility apparatus proper, situating
\cref{thm:historical-distinction} within the trajectory-relative notion
of admissibility that HYDRA uses, and asking what it takes for a system
to allocate history channels in a way that tracks
reconstructive-irreversible transitions specifically, rather than
either everywhere or nowhere.
% ===== END chapters/ch04-two-kinds-of-persistence =====
% ===== BEGIN chapters/ch05-formalizing-irreversibility =====
\chapter{Formalizing Irreversibility}
\label{ch:formalizing-irreversibility}

\cref{ch:two-kinds-of-persistence} established when a history channel
is doing real work: exactly when it is injective on the fiber of
predecessors a transition collapses (\cref{thm:historical-distinction}).
What it did not establish is \emph{which} fibers a system ought to
care about distinguishing in the first place. Not every collapsed
distinction matters. A system that tried to preserve distinguishability
over every non-injective transition it ever executed would drown in
its own history channel long before the retention bought it anything
useful. This chapter supplies the missing criterion, using the notion
--- central to the HYDRA architecture --- that admissibility is a
relational property of a state with respect to a trajectory and a
family of permitted continuations, rather than an intrinsic property
of the state alone.

\section{Admissibility as a Relational Property}

\begin{definition}[Continuation family]
\label{def:continuation-family}
Given a transition system $(X, U, F)$ as in
\cref{def:transition-system}, a continuation family assigns to each
state $x \in X$ a set $C(x) \subseteq X^{\mathbb{N}}$ of trajectories
beginning at $x$ that are considered permitted --- consistent with
whatever constraints (physical, logical, normative, or task-specific)
the system is operating under.
\end{definition}

\begin{definition}[Trajectory-relative admissibility]
\label{def:admissibility}
A state $x_i$ occurring within a trajectory $\tau = (x_0, x_1, \dots)$
is admissible relative to $\tau$ and a continuation family $C$ if the
suffix of $\tau$ beginning at $x_i$ is a member of $C(x_i)$: that is,
if the way the trajectory actually continues after $x_i$ is one of the
continuations $x_i$ permits.
\end{definition}

\begin{remark}
\cref{def:admissibility} makes admissibility relational in the precise
sense HYDRA's architecture requires: it is not a predicate on $x_i$
alone, since the same state can be admissible relative to one
continuation and inadmissible relative to another. This is why
admissibility cannot be checked by inspecting a state in isolation ---
it requires access to the trajectory the state is embedded in, or at
minimum to enough of the trajectory's future to evaluate membership in
$C(x_i)$.
\end{remark}

This relational character is precisely what connects admissibility to
the machinery of \cref{ch:two-kinds-of-persistence}. If admissibility
depended only on the current state, a system could always re-derive it
from $x_i$ and would never need a history channel to evaluate it. Because
admissibility depends on trajectory membership, and because
\cref{prop:fiber-collapse} tells us trajectories are generically
collapsed by $F$, a system's ability to evaluate its own admissibility
--- now or later --- can itself be destroyed by the same collapse that
destroys predecessor identity.

\section{Admissibility Estimates and Repair}
\label{sec:estimates-and-repair}

A terminological note is necessary before proceeding. In the HYDRA
literature, \enquote{projection} has a specific, prior meaning: the
Minimal Projection Theorem constructs, for a fixed admissibility
structure $A$ on a trajectory space $\mathcal{X}$, a unique canonical
quotient map $\pi^* : \mathcal{X} \to \mathcal{M}^*$ collapsing exactly
the states with identical admissible futures ($x_1 \sim x_2 \iff
A(x_1) = A(x_2)$), through which every other admissibility-preserving
map factors. That object is static, canonical, and not something a
system revises --- it is fixed once $A$ is fixed, and "repairing" it
would only make sense as revising one's beliefs about what $A$ is, not
as an operation on $\pi^*$ itself.

The object this chapter needs is different in kind: not the canonical
map $\pi^*$, but a system's own time-indexed, resource-bounded
\emph{estimate} of it, built and corrected under evidence the system
receives incrementally rather than derived in closed form from a
fully specified $A$. We write this object $\hat\pi_t$ throughout, to
keep it visibly distinct from $\pi^*$, and return to the precise
relationship between the two --- $\hat\pi_t$ is best understood as a
sequence converging toward, but not generally reaching,
$\pi^*$ --- in \cref{ch:hydra-projections}.

HYDRA does not evaluate admissibility against the raw state $x_t$
directly. It maintains this intermediate estimate.

\begin{definition}[Admissibility estimate]
\label{def:projection}
An admissibility estimate $\hat\pi$ is a partial model of the
environment maintained by a system: a structured, inspectable
hypothesis about which states are admissible, distinct from $x_t$
itself and revisable independently of it. The space of admissibility
estimates is written $\hat\Pi$.
\end{definition}

\begin{definition}[Fit relation and anomaly]
\label{def:fit}
A fit relation $\vDash \subseteq \hat\Pi \times E$ relates
admissibility estimates to evidence, where $E$ is a space of
observations the system can receive. A piece of evidence $e$ is
anomalous with respect to $\hat\pi$ if $\hat\pi \not\vDash e$: the
estimate fails to accommodate the observation.
\end{definition}

\begin{definition}[Repair operator]
\label{def:repair}
A repair operator is a function $R : \hat\Pi \times E \to \hat\Pi$
such that, for any $\hat\pi$ and any $e$ anomalous with respect to
$\hat\pi$, $R(\hat\pi, e) \vDash e$: the repaired estimate
accommodates the evidence that falsified its predecessor.
\end{definition}

A repair event is a transition in exactly the sense of
\cref{def:transition-system}, with $\hat\Pi$ playing the role of $X$, $E$
playing the role of $U$, and $R$ playing the role of $F$:
\[
\hat\pi_{t+1} = R(\hat\pi_t, e_t).
\]

\begin{proposition}[Repair transitions generically collapse fibers]
\label{prop:repair-collapse}
For a repair operator $R$ satisfying \cref{def:repair}, it is
generically the case --- true for most fit relations $\vDash$ of
practical interest, and false only when $\vDash$ is fine enough to
make every estimate uniquely falsifiable in exactly one way --- that
$R$ is non-injective: distinct pairs $(\hat\pi, e) \neq (\hat\pi', e')$
can satisfy $R(\hat\pi, e) = R(\hat\pi', e')$.
\end{proposition}

\begin{proof}[Justification]
Repair operators are typically required to produce an estimate that
accommodates the new evidence and is otherwise minimal or canonical
in some further sense (simplest correction, least-disruptive revision,
maximum-likelihood update). Whenever more than one prior estimate
$\hat\pi, \hat\pi'$, subjected to more than one piece of anomalous
evidence $e, e'$, converges under this minimality requirement on the
same corrected estimate, $R$ is non-injective at that estimate by
\cref{def:fiber}. This is the ordinary case: many different specific
errors, once corrected under a "fix the anomaly with minimal
disruption" policy, converge on the same repaired state, because the
correction is determined more by the minimality criterion than by the
specific defect that triggered it.
\end{proof}

\cref{prop:repair-collapse} is the bridge \cref{ch:two-kinds-of-persistence}'s
apparatus needed. Repair is not a special kind of transition exempt
from fiber collapse; it is, if anything, an especially prone case,
because minimality criteria actively erase the specific character of
whatever they correct. \cref{thm:historical-distinction} therefore
applies to repair directly: whether the identity of a failed
estimate survives its own repair depends on whether the system's
retained history channel is injective over the fiber of estimates
that repair could have corrected to the same result.

\section{Admissibility-Relevant Fibers}

\cref{cor:retention-not-adequate} already warned that retaining
\emph{some} history channel is not the same as retaining an adequate
one. This section supplies the criterion for adequacy that
\cref{ch:two-kinds-of-persistence} deliberately left open: adequate
\emph{with respect to what}. The answer, given
\cref{def:admissibility}, is: with respect to whichever fibers affect
the evaluation of admissibility for some future state.

\begin{definition}[Admissibility-relevant fiber]
\label{def:admissibility-relevant}
A fiber $\mathrm{Fib}_P(y)$ is admissibility-relevant if there exists
some later state $x_j$, reachable by the system's dynamics from $y$,
and some continuation family $C$, such that the admissibility of $x_j$
relative to $C$ (\cref{def:admissibility}) differs depending on which
member of $\mathrm{Fib}_P(y)$ actually occurred.
\end{definition}

\begin{example}
Suppose a coding agent's context-compaction routine collapses two
distinct reasoning traces --- one that correctly ruled out an unsafe
refactor, and one that merely ran out of budget before considering it
--- into the same summarized context. If a later step re-proposes the
same refactor, whether that proposal is admissible (has it already
been correctly ruled out, or does it merely appear unconsidered)
depends on which predecessor occurred. The compaction fiber is
admissibility-relevant. By contrast, if two different low-level tool
invocations happen to produce byte-identical log lines, and nothing
downstream ever depends on which specific invocation produced the
line, the corresponding fiber is not admissibility-relevant, even
though it is, in the sense of \cref{def:fiber}, collapsed.
\end{example}

\begin{theorem}[Selective Retention]
\label{thm:selective-retention}
A system that wishes to preserve exactly the admissibility judgments
its continuation families can support, at minimum retention cost, must
allocate history channels that are injective
(\cref{def:distinguishability}) precisely over admissibility-relevant
fibers (\cref{def:admissibility-relevant}), and need not allocate
injective history channels over fibers that are not
admissibility-relevant.
\end{theorem}

\begin{proof}[Sketch]
Necessity: if a fiber is admissibility-relevant, there is some future
state $x_j$ and continuation family $C$ whose admissibility verdict
depends on which predecessor in the fiber occurred. If the system's
history channel is not injective over that fiber, then by
\cref{thm:historical-distinction} the predecessor's identity is not
recoverable from $(y, h_{t+1})$, so the admissibility verdict for
$x_j$ cannot be determined either, since by hypothesis it depends on
that identity. Injectivity over the fiber is therefore required to
support the admissibility judgment. Sufficiency of restricting
injective retention to admissibility-relevant fibers, and the minimum-
cost claim, follow from the fact that (by
\cref{def:admissibility-relevant}) no future admissibility judgment
depends on distinguishing predecessors within a non-admissibility-
relevant fiber; retaining injectivity there would preserve a
distinction with no consumer, at nonzero retention cost, and could be
dropped without any admissibility judgment becoming unsupportable.
\end{proof}

\cref{thm:selective-retention} is the theorem that gives the earlier,
informal prediction --- that trace allocation should track historical
non-recoverability rather than mere computational expense --- its
sharpest form. It does not merely say retention should track
irreversibility in general; it says retention should track the
specific subset of reconstructively irreversible fibers whose collapse
would otherwise make some future admissibility judgment
undecidable, and it licenses discarding distinguishability everywhere
else, including at other reconstructively irreversible fibers that
happen not to feed into any later admissibility check.

\section{What This Rules Out}

\cref{thm:selective-retention} has a corollary worth stating on its
own, because it is the form in which the theorem will be tested
against the Claude Code leak in \cref{ch:reading-the-leak}.

\begin{corollary}[Two failure modes]
\label{cor:two-failure-modes}
A system's retention architecture can fail \cref{thm:selective-retention}
in exactly two ways: it can under-retain, leaving some
admissibility-relevant fiber without an injective history channel, so
that some future admissibility judgment becomes undecidable; or it can
over-retain, maintaining injective history channels over fibers that
are not admissibility-relevant, paying a retention cost that no future
admissibility judgment consumes.
\end{corollary}

Under-retention is the failure mode the original \enquote{wizard behind
the curtain} debate never considered, because it is invisible from the
outside until the specific future state that needed the distinction
actually arrives and cannot be evaluated. Over-retention is the
failure mode naive economic accounts
(\cref{ch:why-structure-externalizes}) are well equipped to notice and
penalize, since it shows up directly as wasted storage or compute. The
interesting architectural question --- taken up properly once HYDRA's
own projection machinery is on the table in \cref{ch:hydra-projections}
--- is whether a system's actual pattern of retention tracks
admissibility-relevance specifically, or merely approximates it by
over-retaining broadly (safe but wasteful) or under-retaining broadly
(cheap but occasionally undecidable), neither of which requires the
system to have represented admissibility-relevance as such at all.
% ===== END chapters/ch05-formalizing-irreversibility =====
% ===== BEGIN chapters/ch06-falsifiability =====
\chapter{A Falsifiability Criterion}
\label{ch:falsifiability}

\cref{thm:selective-retention} states what a system's retention
architecture \emph{ought} to track, given the definitions
\cref{ch:formalizing-irreversibility} supplies. Before using it to
read HYDRA's own architecture (\cref{ch:hydra-projections}) or the
Claude Code leak (\cref{ch:reading-the-leak}), it is worth being
honest about what kind of claim it is, because the honest answer
determines what would actually count as evidence against it.

\section{What Would It Mean to Falsify a Theorem Proved From Definitions?}

\cref{thm:selective-retention} was proved, in
\cref{ch:formalizing-irreversibility}, directly from
\cref{def:admissibility-relevant}: a fiber counts as
admissibility-relevant exactly when some future admissibility
judgment depends on distinguishing its members, and the theorem then
observes that injective retention is required exactly there and
nowhere else. This is close to definitional in the same sense that
the HYDRA architecture's own Recursive Admissibility Stabilization
result is close to definitional given its notion of an
$A$-preserving transition: the inductive step follows immediately
once the definitions are fixed, and no substantive empirical claim
about the world has yet been made. A theorem of this kind cannot be
falsified by any observation, because no observation could conflict
with a claim that is true by construction. Its value, as with HYDRA's
result, is organizational: it tells you exactly what property to look
for, not whether that property holds of anything real.

The empirical content of this monograph's argument, if it has any, has
to live elsewhere. Two questions carry that content, and this chapter
addresses both.

\begin{enumerate}
\item Can admissibility-relevance (\cref{def:admissibility-relevant})
be determined for a real fiber in a real system \emph{without}
already knowing what that system retains --- or does the notion
collapse into the same vacuity \cref{ch:why-structure-externalizes}
diagnosed in the naive reconstruction-cost account, where any
retained structure can be redescribed after the fact as having been
worth retaining?
\item If admissibility-relevance can be determined independently,
does a real system's actual pattern of retention track it --- and if
it does not, what observable consequence follows?
\end{enumerate}

\section{An Independence Criterion}

The vacuity risk is real and specific. \cref{ch:why-structure-externalizes}
showed that \enquote{expensive to reconstruct} can be applied to
anything a system happens to retain, after the fact, with no
independent check. \cref{def:admissibility-relevant} is at risk of the
same move: nothing so far stops someone from declaring a fiber
admissibility-relevant precisely because, and only because, the system
under study happens to retain a distinguishing trace there. What is
needed is a way to decide admissibility-relevance that does not
consult the system's retention architecture at all.

\begin{definition}[Externally testable admissibility-relevance]
\label{def:external-relevance}
A fiber $\mathrm{Fib}_P(y)$ is externally admissibility-relevant at a
later state $x_j$, relative to a continuation family $C$, if replaying
the system's dynamics forward from each distinct predecessor $(x,u),
(x',u') \in \mathrm{Fib}_P(y)$ in turn --- holding every subsequent
input fixed --- yields different admissibility verdicts for $x_j$
under $C$ (\cref{def:admissibility}), \emph{regardless of whether the
system's own architecture actually retains anything that would let it
compute this difference}.
\end{definition}

\cref{def:external-relevance} is an interventionist test: it asks what
would have happened under two counterfactual histories, holding
everything downstream fixed, and checks whether the two replays
diverge in their eventual admissibility verdict. It can be carried
out --- in principle, and in practice for systems simple enough to
simulate or log exhaustively --- by an external observer with access
to the system's specification and dynamics, without any reference to
what the system itself chose to keep. This is what makes it usable as
an independence check.

\begin{proposition}[Non-vacuity]
\label{prop:non-vacuity}
Admissibility-relevance, tested via \cref{def:external-relevance}, is
decidable (in the sense of being determinable by simulation or replay
against the system's specified dynamics) without reference to the
system's retention architecture. Consequently, the prediction of
\cref{thm:selective-retention} --- that retention should track
admissibility-relevance --- is not vacuous in the way
\cref{ch:why-structure-externalizes} showed the naive reconstruction-
cost account to be: relevance and retention are established by two
independent procedures, and it is possible, prior to inspecting a
system's retention architecture, to have already fixed which of its
fibers are relevant.
\end{proposition}

This is the qualification that matters: \cref{prop:non-vacuity} does
not claim any \emph{particular} system's retention tracks relevance.
It claims only that the question is well-posed --- that there is a
fact of the matter, checkable independently, against which a system's
actual retention pattern can be compared and found to match or not to
match.

\section{The Audit Procedure}

Put together, \cref{def:external-relevance} and
\cref{thm:selective-retention} license a concrete procedure,
applied to HYDRA's own architecture in \cref{ch:hydra-projections} and
to the Claude Code leak in \cref{ch:reading-the-leak}.

\begin{enumerate}
\item \textbf{Enumerate candidate fibers.} Inspect the system's
transition dynamics for non-injective transitions
(\cref{def:fiber}) --- edits, corrections, compactions, merges,
overwrites, and any other operation that is many-to-one on admissible
predecessors.
\item \textbf{Test relevance externally.} For each candidate fiber,
apply \cref{def:external-relevance}: replay from distinct
predecessors with subsequent inputs held fixed, and check whether any
downstream admissibility verdict diverges.
\item \textbf{Audit retention.} Independently, inspect what the
system actually retains at that fiber, and check whether the retained
structure is injective on the fiber (\cref{def:distinguishability}).
\item \textbf{Compare.} Cross-tabulate the two independent findings
against \cref{cor:two-failure-modes}: relevant and retained is
compliant; relevant and not retained is under-retention; not relevant
and retained is over-retention; not relevant and not retained is
compliant by omission.
\end{enumerate}

\section{Two Failure Signatures}

The comparison in step 4 produces a prediction with observable
consequences, and it is worth stating what those consequences should
look like, because \enquote{under-retention} and \enquote{over-
retention} are not directly visible from the audit alone --- they are
claims about what a system's behavior \emph{should} do, which still
needs to be checked against what it \emph{does} do.

\begin{definition}[Undecidability failure]
\label{def:undecidability-failure}
A system exhibits an undecidability failure at $x_j$ if a downstream
admissibility judgment for $x_j$ is, by \cref{def:external-relevance},
genuinely relevance-dependent on some fiber, the system's retention
does not preserve distinguishability over that fiber
(\cref{def:distinguishability}), and the system nonetheless produces a
verdict for $x_j$ --- necessarily either by guessing, by falling back
to a default, or by silently picking one member of the fiber as if it
were the only one.
\end{definition}

\cref{def:undecidability-failure} is stated carefully because the
naive expectation --- that under-retention shows up as a visible
crash or error --- is usually wrong. Systems under-provisioned in this
sense do not typically halt; they resolve the ambiguity silently,
often by defaulting to whichever fiber member is most recent, most
common, or otherwise privileged by the implementation, and proceed as
if no ambiguity existed. This makes the failure mode empirically
subtler than a crash, but not empirically inaccessible.

\begin{definition}[Default-resolution test]
\label{def:default-resolution-test}
Where an undecidability failure is suspected, construct a controlled
pair of runs that differ only in which member of the relevant fiber
occurred, with everything else held fixed, and check whether the
system's eventual verdict for $x_j$ tracks the true predecessor across
repeated trials at a rate exceeding chance. If the verdict is
statistically independent of which predecessor actually occurred, the
system is resolving the ambiguity by default rather than by genuine
distinction, confirming an undecidability failure by
\cref{def:undecidability-failure} even though no explicit error was
raised.
\end{definition}

Over-retention is comparatively easy to confirm: it shows up directly
as retained structure, chargeable against storage or compute, that no
downstream admissibility judgment --- checked exhaustively via
\cref{def:external-relevance} against every state reachable from the
fiber in question --- ever consults.

\section{Two Falsification Conditions}

With \cref{def:undecidability-failure} and
\cref{def:default-resolution-test} in hand, the empirical claim behind
\cref{thm:selective-retention} --- that a well-engineered system's
retention pattern should approximate the relevance/retention match
\cref{cor:two-failure-modes} describes, and that departures from the
match should correlate with observable consequences --- can be stated
in a form that a specific system can fail.

\begin{itemize}
\item \textbf{Falsification by silent success.} If a system is found,
by \cref{def:external-relevance}, to under-retain at a
genuinely admissibility-relevant fiber, and the default-resolution
test (\cref{def:default-resolution-test}) shows its verdicts at the
affected downstream states nonetheless track the true predecessor at
better than chance --- meaning some mechanism other than the audited
retention channel is carrying the distinction --- the audit procedure
has misattributed the system's competence, and the specific claim that
\emph{this} retention channel is what supports \emph{this}
admissibility judgment is false. This does not refute
\cref{thm:selective-retention} as a definitional claim, but it refutes
the specific reading of a system as an instance of the theorem's
prescription, and would require locating the channel that is actually
doing the work.
\item \textbf{Falsification by uncorrelated retention.} If a system's
overall retention pattern, audited across many fibers, shows no
tendency for admissibility-relevant fibers (per
\cref{def:external-relevance}) to receive injective treatment more
often than admissibility-irrelevant fibers do --- if relevance and
retention are statistically independent across the audited fibers ---
then the broader claim that engineering pressure shapes retention
architecture toward admissibility-relevance, rather than toward
something else entirely (raw reconstruction cost, implementation
convenience, historical accident), is falsified for that system. This
is the condition \cref{ch:reading-the-leak} will actually test.
\end{itemize}

\section{Scope Limits}

Two limitations of the audit procedure are worth stating plainly
rather than discovering by way of an overreached claim later.

First, exhaustive enumeration of fibers (step 1 of the audit
procedure) is generally infeasible for any system complex enough to be
interesting; real audits will sample candidate fibers rather than
enumerate them, and a sampled audit can only support probabilistic
conclusions about the overall correlation claim, not certainty about
any specific unsampled fiber.

Second, \cref{def:external-relevance}'s counterfactual replay assumes
the system's dynamics are specified precisely enough to replay at all.
For a system with genuinely stochastic or environment-coupled
transitions, \enquote{holding subsequent inputs fixed} requires care
about what counts as the same subsequent input across two runs that
started from different predecessors, and the test degrades gracefully
into a statistical rather than a deterministic comparison. This is not
a defect specific to this framework; it is the ordinary difficulty of
counterfactual evaluation in any non-deterministic system, and the
default-resolution test (\cref{def:default-resolution-test}) is
already built to accommodate it by using repeated trials rather than a
single comparison.

Neither limitation undermines the two falsification conditions above;
both mean that a real audit produces a graded, evidence-weighted
verdict rather than a single decisive proof. \cref{ch:hydra-projections}
applies the procedure to HYDRA's own projection machinery, where the
architecture is specified precisely enough for a fairly direct audit.
\cref{ch:reading-the-leak} applies it to the Claude Code leak, where
the audit is necessarily sampled and partial, and the chapter says so
explicitly rather than overstating what a leaked codebase, inspected
after the fact, can actually establish.
% ===== END chapters/ch06-falsifiability =====

\part{HYDRA and the Case Study}
% ===== BEGIN chapters/ch07-hydra-projections =====
\chapter{HYDRA's Projections as Explicit Intermediate Representation}
\label{ch:hydra-projections}

\cref{ch:formalizing-irreversibility} left one question open by
design: whether iterating the repair operator $R$ actually drives the
estimate sequence $\hat\pi_0, \hat\pi_1, \hat\pi_2, \dots$ toward the
canonical quotient $\pi^*$ that HYDRA's Minimal Projection Theorem
guarantees exists, or merely toward \emph{some} sequence that keeps
accommodating whatever evidence arrives without ever becoming correct.
This chapter answers the question, and the answer is more interesting
than the affirmative most treatments would reach for: consistency,
stability, and convergence are three different properties, the repair
operator as defined so far only guarantees the weakest of them, and
the gap between stability and convergence has an exact, constructible
counterexample rather than remaining a vague worry.

\section{Repair as Hypothesis Pruning}

The cleanest way to state the distinction is to stop treating
$\hat\pi_t$ as a free-floating \enquote{estimate} and instead say
explicitly what it is an estimate \emph{of}.

\begin{definition}[Hypothesis space]
\label{def:hypothesis-space}
Let $\mathcal{A}$ be the space of admissibility structures a system
considers possible on its trajectory space $\mathcal{X}$ --- the
candidates among which the true structure $A^* \in \mathcal{A}$ lies.
Each $A \in \mathcal{A}$ induces, by the Minimal Projection Theorem, a
canonical quotient $\pi_A : \mathcal{X} \to \mathcal{M}_A$; write
$\pi^* := \pi_{A^*}$ for the canonical quotient induced by the true
structure.
\end{definition}

\begin{definition}[Estimate as hypothesis set]
\label{def:estimate-as-set}
A system's admissibility estimate at time $t$ is the set $H_t
\subseteq \mathcal{A}$ of admissibility structures still compatible
with everything accommodated so far. The induced estimate quotient is
$\hat\pi_t := \pi_{H_t}$, defined by $x_1 \sim_{H_t} x_2$ iff $\pi_A(x_1)
= \pi_A(x_2)$ for every $A \in H_t$: two states are treated as
equivalent by the estimate only when every remaining candidate
structure agrees they are equivalent.
\end{definition}

\cref{def:estimate-as-set} makes the estimate conservative by
construction: it only asserts an identification when no live
hypothesis would contest it. This is the right default for an
object that is supposed to support admissibility judgments
(\cref{def:admissibility}) rather than merely summarize evidence,
since a false identification --- treating two states as equivalent
when the true structure distinguishes them --- is a different and
worse kind of error than a missing identification.

\begin{definition}[Repair as pruning]
\label{def:repair-pruning}
A repair operator $R$ acting on hypothesis sets satisfies:
\begin{enumerate}
\item \emph{Fit.} Every $A \in R(H_t, e_t)$ is compatible with $e_t$
(the evidence that triggered repair).
\item \emph{Contraction.} $R(H_t, e_t) \subseteq H_t$: repair only
removes candidates, never reintroduces a structure previously
excluded. This is the single axiom subsuming what earlier discussion
treated as two separate conditions --- distinguishability monotonicity
and non-reintroduction of eliminated equivalence classes are the same
requirement stated on the estimate side and the hypothesis side of the
same object.
\item \emph{Soundness.} If the evidence stream $e_0, e_1, \dots$ is
genuinely generated by $A^*$, then $A^* \in H_t$ for all $t$: repair
never prunes the true structure, because true evidence never
contradicts it.
\end{enumerate}
\end{definition}

Fit is the condition \cref{def:repair} already stated. Contraction and
Soundness are the additions this chapter makes explicit, and together
they replace the earlier, weaker requirement --- that a repaired
estimate merely accommodate the \emph{latest} anomaly --- with the
stronger requirement that it remain compatible with the \emph{entire}
evidence history. This is exactly the gap identified against
\cref{def:repair} as originally stated: nothing there prevented a
repair from silently reintroducing an inconsistency with $e_{t-1}$
while fixing its inconsistency with $e_t$. \cref{def:repair-pruning}
closes that gap by axiom rather than leaving it to be assumed.

\section{Consistency, Stability, and Convergence}

With $H_t$ as the primitive object, the three notions can be stated
without ambiguity.

\begin{definition}[Consistency]
\label{def:consistency}
The estimate sequence is consistent if $H_t \neq \emptyset$ for all
$t$. By Soundness (\cref{def:repair-pruning}), consistency is
automatic whenever the evidence is genuinely generated by some fixed
$A^* \in \mathcal{A}$: $A^*$ itself always witnesses non-emptiness.
\end{definition}

\begin{proposition}[Stability under a descending chain condition]
\label{prop:stability}
If $\mathcal{A}$ satisfies the descending chain condition (every
strictly decreasing chain of subsets of $\mathcal{A}$ is finite), then
the sequence $H_0 \supseteq H_1 \supseteq H_2 \supseteq \cdots$
guaranteed by Contraction (\cref{def:repair-pruning}) stabilizes: there
exists $T$ such that $H_t = H_T$ for all $t \geq T$. Write $H_\infty :=
H_T$ for this eventual value.
\end{proposition}

\begin{proof}
A strictly decreasing chain $H_0 \supsetneq H_1 \supsetneq H_2
\supsetneq \cdots$ would violate the descending chain condition unless
it terminates; since $H_t \supseteq H_{t+1}$ always by Contraction, the
sequence either strictly decreases (finitely often, by hypothesis) or
stabilizes, and once $H_t = H_{t+1}$, Contraction forces $H_{t'} = H_t$
for all $t' \geq t$.
\end{proof}

\cref{prop:stability} is doing exactly the work anticipated earlier:
stability is not a further axiom on $R$, it is a free consequence of
Contraction plus a well-foundedness assumption on the hypothesis space.
Systems with infinite, non-well-founded hypothesis spaces need not
stabilize at all; this is a genuine restriction, not a technicality,
and is returned to in \cref{sec:hydra-scope}.

\begin{definition}[Convergence]
\label{def:convergence}
The estimate sequence converges (to $\pi^*$) if $H_\infty$ induces the
same quotient as $A^*$: $\pi_{H_\infty} = \pi^*$.
\end{definition}

\begin{proposition}[The limit estimate refines the truth]
\label{prop:limit-refines}
$\pi_{H_\infty}$ is always at least as fine as $\pi^*$: whenever
$\pi_{H_\infty}(x_1) = \pi_{H_\infty}(x_2)$, also $\pi^*(x_1) =
\pi^*(x_2)$. Equivalently, $\pi_{H_\infty}$ never falsely identifies
two states that $\pi^*$ distinguishes.
\end{proposition}

\begin{proof}
By \cref{def:estimate-as-set}, $\pi_{H_\infty}(x_1) =
\pi_{H_\infty}(x_2)$ requires $\pi_A(x_1) = \pi_A(x_2)$ for every $A
\in H_\infty$. By Soundness, $A^* \in H_\infty$, so in particular
$\pi_{A^*}(x_1) = \pi_{A^*}(x_2)$, i.e., $\pi^*(x_1) = \pi^*(x_2)$.
\end{proof}

\cref{prop:limit-refines} is the precise sense in which the estimate
is conservative rather than merely cautious: it can fail to converge
only by retaining \emph{too many} distinctions, never by collapsing
ones the truth requires. This gives an exact characterization of when
convergence holds and, more usefully, an exact description of what a
failure to converge looks like.

\begin{theorem}[Convergence characterization]
\label{thm:convergence-characterization}
The estimate sequence converges to $\pi^*$ (\cref{def:convergence}) if
and only if every $A \in H_\infty$ induces the same quotient as $A^*$
--- equivalently, if and only if the evidence stream eventually
separates $A^*$, in the sense of inducing a different quotient, from
every rival hypothesis in $\mathcal{A}$ that disagrees with $A^*$ on
the classification of at least one pair of states.
\end{theorem}

\begin{proof}
By \cref{prop:limit-refines}, $\pi_{H_\infty}$ never identifies a pair
$\pi^*$ distinguishes, so $\pi_{H_\infty} = \pi^*$ can fail only by
$\pi_{H_\infty}$ withholding an identification $\pi^*$ makes. By
\cref{def:estimate-as-set}, $\pi_{H_\infty}$ withholds the
identification of $x_1, x_2$ exactly when some $A' \in H_\infty$
dissents, i.e., $\pi_{A'}(x_1) \neq \pi_{A'}(x_2)$ while
$\pi^*(x_1) = \pi^*(x_2)$ --- which is precisely a pair on which $A'$
disagrees with $A^*$. So $\pi_{H_\infty} = \pi^*$ holds iff no
surviving $A' \in H_\infty$ disagrees with $A^*$ on any pair, i.e.,
iff every surviving hypothesis induces the identical quotient to
$A^*$.
\end{proof}

\begin{example}[Stable, consistent, and non-convergent]
\label{ex:non-convergence}
Let $\mathcal{X} = \{x_1, x_2, x_3\}$ and let $\mathcal{A} = \{A^*,
A'\}$ where $A^*$ places $x_1$ and $x_2$ in the \emph{same}
admissibility class, $A'$ places them in \emph{distinct} classes, and
both structures agree on every other classification, including
everything about $x_3$. Suppose the evidence stream $e_0, e_1, \dots$
consists entirely of observations bearing on $x_3$ and on the shared
classifications, never on whether $x_1$ and $x_2$ are
admissibility-equivalent. Then $A^* \in H_t$ for all $t$ by Soundness,
and $A'$ is never excluded either, since no evidence has ever borne on
the one point where $A^*$ and $A'$ disagree; so $H_t = \{A^*, A'\}$
for all $t$, the sequence is trivially stable from $t = 0$, and it is
consistent throughout. But $\pi^*(x_1) = \pi^*(x_2)$, while
$\pi_{H_t}(x_1) \neq \pi_{H_t}(x_2)$ for every $t$, since $A'$'s
dissent blocks the identification by \cref{def:estimate-as-set}. The
sequence is stable and consistent for all $t$, and never converges: it
permanently withholds an identification the truth supports, because
the one piece of evidence that would resolve $A^*$ against $A'$ never
arrives.
\end{example}

\cref{ex:non-convergence} is the constructible version of \enquote{a
projection can be stable without being correct}: stability is a
property of the evidence stream drying up on the disagreements that
remain, not a property of the disagreements having been resolved.
Convergence requires the stronger condition that the evidence stream
be, in the relevant sense, exhaustive with respect to every rival
hypothesis --- a condition \cref{thm:convergence-characterization}
states exactly and \cref{ex:non-convergence} shows is not automatic.

\section{Repair Logs as Admissibility-Relevant History Channels}

The hypothesis-pruning formulation also answers a question
\cref{ch:formalizing-irreversibility} left implicit: what, concretely,
should a system retain around a repair event, beyond the repaired
estimate itself? The natural candidate is now visible directly from
\cref{def:repair-pruning}: the record of which hypotheses were pruned,
and by which evidence.

\begin{definition}[Repair log]
\label{def:repair-log}
A repair log at time $t$ is a record of the pruning step $H_{t-1}
\setminus H_t$ together with the evidence $e_{t-1}$ responsible for it.
\end{definition}

By \cref{thm:selective-retention}, this log is admissibility-relevant
--- worth retaining as an injective history channel in the sense of
\cref{def:distinguishability} --- exactly when some later comparison
between surviving hypotheses in $H_{t'}$, $t' > t$, would come out
differently depending on which specific structures were pruned at
step $t$ rather than merely how many. This is not automatic: two
different pruning steps that happen to leave $H_t$ in the same state
are, by \cref{prop:fiber-collapse}, indistinguishable from $H_t$
alone, and \cref{thm:selective-retention} says the log is worth
keeping only insofar as some downstream admissibility judgment
actually depends on knowing which. A repair log kept unconditionally,
for every pruning step regardless of downstream relevance, is exactly
the over-retention failure mode of \cref{cor:two-failure-modes}.

\section{What Makes This an Explicit Intermediate Representation}
\label{sec:explicit-intermediate}

\cref{ch:category-mistake} left open what it would take for an
artificial system to possess something functionally comparable to
Logical Form: an intermediate representation, distinct from both raw
input and final interpretation, that structure-building operations
construct and later stages consume. The apparatus of this chapter is
offered as a candidate answer, and its distinguishing property is not
that it exists --- interpretability researchers have long since found
structure of various kinds inside transformer latent spaces --- but
that it is \emph{inspectable in the specific sense the audit procedure
of \cref{ch:falsifiability} requires}.

Given a system built on \cref{def:estimate-as-set}, the audit
procedure's steps are not merely available in principle; they are
answerable directly from the architecture's own bookkeeping. Step 1
(enumerate candidate fibers) corresponds to enumerating repair events;
step 3 (audit retention) corresponds to checking which repair logs
(\cref{def:repair-log}) were kept; and \cref{def:external-relevance}'s
counterfactual replay, which for an opaque system requires simulating
the entire architecture twice, here reduces to comparing $H_t$ across
the two hypothetical evidence orderings directly, since $H_t$ is
already an explicit set rather than a distributed pattern that must be
reconstructed by probing. A latent vector must be interpreted before
it can be compared; a hypothesis set can simply be compared.

This is a narrower claim than \enquote{HYDRA understands admissibility
and transformers do not,} and it should be read as exactly that
narrow. What \cref{sec:explicit-intermediate} establishes is that
\emph{if} a system is built on \cref{def:estimate-as-set}, the
question of whether its retention tracks admissibility-relevance is
decidable by inspection rather than by inference. Whether any given
implementation of HYDRA is in fact built this way, and whether
anything comparable can be recovered from a transformer's internals
well enough to run the same audit, are separate empirical questions,
taken up for the leaked architecture specifically in
\cref{ch:reading-the-leak} and for the general transformer semantics
question in \cref{ch:transformer-semantics}.

\section{Scope: What This Chapter Assumes About HYDRA}
\label{sec:hydra-scope}

Two assumptions this chapter's results depend on are worth stating
plainly, because HYDRA's own presentation is honest about exactly
where its formal machinery is established versus conjectural, and this
chapter's results should be filed the same way.

First, \cref{prop:stability} requires $\mathcal{A}$ to satisfy a
descending chain condition. HYDRA's own account of admissibility
structures does not establish this in general; it is a substantive
assumption about the hypothesis space a given implementation searches
over, true for finite or well-ordered hypothesis spaces and not
guaranteed otherwise. Where it fails, an estimate sequence may
continue revising indefinitely without ever stabilizing, and neither
\cref{def:convergence} nor \cref{thm:convergence-characterization}
directly applies.

Second, this chapter's entire apparatus presupposes that a system's
estimate is well-modeled as a subset $H_t$ of some fixed hypothesis
space $\mathcal{A}$, known in advance. Real systems may instead
construct or extend $\mathcal{A}$ itself as evidence arrives --- a
form of ontological revision, rather than mere pruning of a fixed
space, that \cref{def:repair-pruning} does not cover. Whether HYDRA's
actual repair machinery operates over a fixed or an expanding
hypothesis space is left open here and is relevant to how much of
\cref{thm:convergence-characterization} survives contact with a real
implementation --- a question \cref{ch:reading-the-leak} does not
resolve either, since a leaked codebase's implementation choices do
not, by themselves, settle a question about the mathematical structure
of its hypothesis space.
% ===== END chapters/ch07-hydra-projections =====
% ===== BEGIN chapters/ch08-reading-the-leak =====
\chapter{Reading the Leak Through the Distinction}
\label{ch:reading-the-leak}

This chapter applies the audit procedure of \cref{ch:falsifiability} to
the Claude Code leak described in \cref{ch:introduction}. It does not
revisit whether the presence of a parser inside the leaked codebase
says anything about semantic competence; \cref{ch:category-mistake}
settled that. The question here is narrower and more useful: does the
leaked architecture's pattern of externalized structure actually sort
into reconstruction-cost and irreversibility materialization the way
\cref{ch:two-kinds-of-persistence} and \cref{ch:formalizing-irreversibility}
predict, and where it retains history channels around collapsed
fibers, does that retention track admissibility-relevance
(\cref{def:admissibility-relevant}) rather than something else?

\section{What Kind of Audit This Is}
\label{sec:audit-caveat}

\cref{ch:falsifiability} was explicit that the audit procedure
degrades gracefully but not silently when full replay access is
unavailable. This chapter's version of the audit is a
\emph{specification-level} audit: it works from architectural details
made public through post-leak technical analysis --- blog write-ups,
security research, and independent reconstructions of the leaked
TypeScript source --- rather than from direct execution of the
codebase against controlled counterfactual histories. This means
step 2 of the audit procedure, the counterfactual replay of
\cref{def:external-relevance}, is applied here in a weakened form:
rather than literally running the system twice from divergent
predecessors, admissibility-relevance is inferred from documented
behavior --- what the architecture's own hooks, checkpoints, and
permission logic are described as doing --- which is a real but
strictly weaker form of evidence than a controlled replay. Every
classification below should be read with that qualification attached,
and \cref{sec:leak-audit-limits} returns to it directly.

\section{The Audit Table}

\begin{table}[htbp]
\centering
\footnotesize
\setlength{\tabcolsep}{3pt}
\begin{tabular}{p{2.7cm}p{1.4cm}p{2.3cm}p{2.3cm}p{2.1cm}}
\toprule
\textbf{Structure} & \textbf{Generating input?} & \textbf{Admissibility-relevant (documented)?} & \textbf{Retention observed?} & \textbf{Classification} \\
\midrule
Parse/search infrastructure (file indexing, grep-style tools) & Yes --- source tree on disk & No & Cached, not logged & Reconstruction-cost \\
Prompt cache (cache-break vectors) & Yes --- prior context tokens & No & Cached with explicit invalidation conditions & Reconstruction-cost \\
Companion PRNG state (deterministic per-user generator) & Yes --- user identifier plus fixed salt & No & Recomputed, not stored & Reconstruction-cost (trivial) \\
Bash command risk classification (numbered security checks) & N/A --- stateless function of command content & No fiber collapse to track & None needed & Not a retention case \\
Compaction chain (non-destructive boundary markers on the transcript) & Yes, \emph{while} source transcript retained --- see \cref{sec:compaction-ambiguous} & Yes --- later steps must be able to tell what was summarized versus discarded & Retained; boundaries patched at read time; transcript retention duration undocumented & Time-dependent: reconstruction-cost while transcript survives, reclassifies on pruning \\
Pre-risky-operation checkpoint (version-control commit before a flagged mutation) & Yes, but only \emph{because} explicitly created for this purpose & Yes --- rollback admissibility depends on the pre-operation state & Retained, deliberately, keyed to risk classification & Irreversibility \\
Cross-session trust state & No & Contested --- see \cref{sec:session-trust} & Deliberately not retained & See \cref{sec:session-trust} \\
Background daemon decision log (unconfirmed feature) & No & Plausible, unconfirmed & Reportedly append-only & Irreversibility (low confidence) \\
\bottomrule
\end{tabular}
\caption{Specification-level audit of documented Claude Code components against \cref{def:reconstruction-cost} and \cref{def:irreversibility-materialization}.}
\label{tab:leak-audit}
\end{table}

\section{A Confirming Case}

One row of \cref{tab:leak-audit} is worth dwelling on as a clean
confirmation, and it is worth being honest that it is the only one
that stays clean under scrutiny. \cref{sec:compaction-ambiguous}
returns to the compaction chain, which looked equally clean on first
pass and turns out not to be once its generating input is examined
over time rather than at a single instant.

The pre-risky-operation checkpoint confirms the \emph{selective} part
of \cref{thm:selective-retention}. A version-control commit is not
retained before every operation; it is documented as triggered
specifically before operations flagged as risky. This is close to a
direct implementation of the fiber-injectivity test from
\cref{ch:two-kinds-of-persistence}: the architecture is not paying
retention cost uniformly across all transitions, which
\cref{cor:two-failure-modes} would flag as over-retention if it were,
nor is it declining to retain around destructive edits generally,
which would be under-retention; it is retaining precisely at the
transitions where a later admissibility judgment --- is rollback to
the pre-operation state still available and consistent --- would
otherwise become undecidable in the sense of
\cref{def:undecidability-failure}. Unlike the compaction chain, the
checkpoint's status does not depend on anything decaying later: a
git commit, once made, does not become a worse generating input over
time in the way \cref{sec:compaction-ambiguous} shows a compacted
transcript can.

\section{An Ambiguous Case: Context Compaction and the Decay of Generating Inputs}
\label{sec:compaction-ambiguous}

The compaction chain was classified, above and in \cref{tab:leak-audit},
as a clean instance of irreversibility materialization: boundaries
recorded explicitly, original material preserved append-only rather
than destructively edited. That classification is correct for a
system in the state the documentation describes --- but it is correct
only \emph{conditionally}, and making the condition explicit exposes a
gap in \cref{def:generating-input} itself that the earlier chapters did
not need to confront and this case forces into the open.

\cref{def:generating-input} asks whether a state $s$ exists,
independent of the system's own execution history, from which a
structure $X$ can be reproduced. Nothing in that definition says
\emph{when} $s$ must exist. Every use of the definition so far has
been able to get away with this omission because the cases considered
--- source files on disk, corpora underlying a search index --- are
generating inputs that, for practical purposes, do not expire.
Transcripts are different. They are exactly the kind of state that
real systems prune: documented session-lifecycle hooks include cleanup
tasks that remove stale logs and temporary files, and no system
retains raw conversational history indefinitely as storage and cost
considerations accumulate. This forces a question
\cref{def:generating-input} was never asked before: a generating input
for what, \emph{measured at which time}?

\begin{definition}[Time-relative generating input]
\label{def:time-relative-generating-input}
A state $s$ is a generating input for $X$ at time $t$
(\cref{def:generating-input}) if $s$ remains accessible to the system,
or to an auditor applying \cref{ch:falsifiability}'s procedure, at $t$.
Write $\mathrm{GenInput}(X, t) = 1$ if such an $s$ is accessible at
$t$, and $0$ otherwise.
\end{definition}

\begin{proposition}[Classification can decay]
\label{prop:classification-decay}
A structure $X$ can satisfy $\mathrm{GenInput}(X, t_1) = 1$ and
$\mathrm{GenInput}(X, t_2) = 0$ for $t_2 > t_1$, without any change to
$X$ itself, to the process that produced it, or to anything executed
at $t_2$: it suffices that the generating input $s$, accessible at
$t_1$, becomes inaccessible by $t_2$ through an unrelated retention
decision --- a cleanup routine, an archival policy, a storage limit
--- that has nothing to do with $X$ or with the transition that
produced it.
\end{proposition}

\begin{proof}
Immediate from \cref{def:time-relative-generating-input}: the
predicate $\mathrm{GenInput}(X, t)$ depends only on whether $s$ is
accessible at $t$, which can change due to any process that removes
$s$, regardless of whether that process has any relationship to $X$.
\end{proof}

Applied to the compaction chain, \cref{prop:classification-decay} says
something the earlier confirming treatment elided. So long as the
pre-compaction transcript is retained, the compacted summary is, in
the relevant sense, reconstruction-cost material relative to that
transcript: a higher-fidelity re-derivation remains possible, and the
transcript itself is the object actually doing the irreversibility-
grounds work, exactly as the earlier classification assumed. But if a
session's raw transcript is later pruned --- by exactly the kind of
cleanup task documented elsewhere in the same architecture --- every
summary derived from it silently changes classification at that
moment, from a reconstruction-cost derivative of a retained transcript
to the \emph{only} remaining trace of a now-permanently-collapsed
fiber. And because compaction is lossy by design --- it exists
specifically to discard detail under context pressure --- the summary
that remains after that moment is not injective over the fiber it
would now need to preserve (\cref{def:distinguishability}): it was
never built to be, because at the time it was built, it did not need
to be.

This is a genuinely uncomfortable finding for the classification
scheme this monograph has been applying, not because it shows the
theory wrong, but because it shows the theory's own machinery
identifying a failure mode invisible to a single-timepoint audit. The
available reporting does not establish whether the compaction system
tracks, or is designed to care about, whether a summary's source
transcript has since been pruned; there is no documented mechanism
analogous to the pre-risky-operation checkpoint's explicit triggering
condition that would suggest the architecture treats this transition
--- from reconstruction-cost to reconstructive irreversibility --- as
one requiring any response. If no such mechanism exists, the
architecture is, at the moment of pruning, silently converting a
population of compliant structures into under-retained ones without
any audit trail marking that it happened, which is precisely the kind
of failure \cref{def:undecidability-failure} describes and precisely
the kind step 3 of the audit procedure (\cref{ch:falsifiability}) is
least equipped to catch, since it audits retention as it currently
stands rather than tracking how retention adequacy changes as
generating inputs elsewhere expire.

\paragraph{The same pattern recurs twice more.} Two other rows in
\cref{tab:leak-audit} exhibit the identical structure once looked at
this way, and neither is worked through in full here, since the
mechanism is now established. A trust score derived from a history of
approvals and denials is reconstruction-cost material relative to that
history while the history survives; if the history is later discarded
and only the aggregate score retained, the same silent reclassification
occurs, and the score --- a single scalar --- is certainly not
injective over the fiber of specific approval histories that could
have produced it. Repair logs (\cref{def:repair-log}) themselves are
not exempt: \cref{ch:hydra-projections} treated a repair log as the
history channel that preserves distinguishability over a pruning
event, but a repair log is itself a structure subject to its own
later transitions --- archival, compression, summarization under
storage pressure --- and \cref{prop:repair-collapse}'s argument that
repair transitions generically collapse fibers applies recursively to
the log's own lifecycle exactly as it applies to the estimate the log
was built to explain. A compressed repair log is a history channel
that has itself become subject to reconstructive irreversibility,
and nothing this monograph has established guarantees that the
compression preserves distinguishability over whichever fibers the
original log was retained to cover.

\paragraph{A second, sharper tension, independent of pruning.} The
decay argument above concerns what happens to a summary's
classification \emph{over time}, as its generating input expires. A
different and harder problem is already present at the moment the
summary is first created, before any transcript has been pruned at
all. A single compaction event does not act on one fiber; it acts
across many simultaneously --- every distinct fact, decision, and
ruled-out alternative in the pre-compaction transcript is a candidate
fiber in its own right, and the summarizer's output is injective over
some of these and non-injective over the rest, by construction, in the
same breath. This is not a defect to be engineered away: a summary
that was injective over every fiber in the original transcript would
simply be the transcript, and compaction would have accomplished
nothing.

\begin{definition}[Selective compression]
\label{def:selective-compression}
Let $F_1, \dots, F_n$ be the candidate fibers present in a structure
$X$ prior to a compressive transition $G$. Write $S_G \subseteq \{1,
\dots, n\}$ for the subset of fibers over which $G$ is injective
(\cref{def:distinguishability}). $G$ is admissibility-aligned if $S_G$
contains every admissibility-relevant fiber
(\cref{def:admissibility-relevant}) among $F_1, \dots, F_n$;
admissibility-misaligned if some admissibility-relevant fiber falls
outside $S_G$.
\end{definition}

\cref{def:selective-compression} is the sharper version of the
question this case actually poses. It is not \enquote{does the
compacted summary preserve distinguishability,} which has the trivial
answer \emph{partially, obviously, that is what compression means.}
It is \emph{whether the specific partiality lines up with what future
admissibility judgments will need} --- whether $S_G$ happens to cover
the admissibility-relevant fibers specifically, or covers some other
selection entirely, chosen by whatever criterion the summarizer
actually optimizes for. \cref{thm:selective-retention}, applied at
this resolution, is a design target for the summarizer, not merely for
the surrounding architecture: a compaction mechanism that is
admissibility-aligned in the sense of \cref{def:selective-compression}
is doing exactly what \cref{thm:selective-retention} would prescribe
at the level of individual facts, and one that is misaligned is
producing exactly the under-retention failure of
\cref{cor:two-failure-modes}, fact by fact, regardless of how
sophisticated or well-reviewed its overall summarization quality
appears to be by any other standard.

Nothing in the available reporting settles whether Claude Code's
compaction is admissibility-aligned in this sense, and it is worth
being explicit about why this is a harder question than anything else
audited in this chapter. The parser, the checkpoint, and the trust
score are all governed by code whose retention behavior is, at least
in principle, inspectable directly: a reviewer can read the triggering
condition. A compaction step performed by the model itself has no
comparable artifact to inspect --- \enquote{what did the summarizer
decide was worth keeping, and why} is not a policy written down
anywhere; it is a judgment made by the same kind of computation
\cref{ch:transformer-semantics} already flagged as functionally
adequate in ways that are not always architecturally transparent
(\cref{def:functional-adequacy,def:architectural-transparency}). This
means the compaction case is not merely unresolved by this chapter's
specification-level audit; it may be a case where resolving it
requires exactly the protocol \cref{ch:transformer-semantics} proposed
and did not carry out --- constructing disambiguating contexts,
compacting them, and checking whether the admissibility-relevant fact
survives the compaction --- applied to the summarizer itself rather
than to the base model's ordinary forward pass. The audit procedure
does not fail silently here; it identifies, correctly, that the
question it was built to answer has, in this one case, been handed to
a component this monograph has already argued cannot be audited the
same way everything else in this chapter was.

None of this identifies a case satisfying the strongest disconfirming
pattern \cref{sec:leak-audit-falsification} describes --- a structure
with no generating input, a non-injective retained channel, and
behavior that nonetheless tracks the true predecessor, indicating some
other mechanism is quietly carrying the distinction. Finding that case
would require exactly the replay-level access
\cref{sec:leak-audit-limits} already flagged this chapter as lacking:
the ability to test, against the running system rather than its
documentation, whether a post-pruning admissibility judgment that
\emph{should} be undecidable by this chapter's own analysis is in fact
answered correctly, and if so, by what. What this section establishes
is narrower and, on its own terms, still a real complication: that
the classification scheme has a time axis this monograph did not
originally give it, that it also has a within-instant resolution axis
this monograph did not originally give it either, that at least one
leaked component sits exactly on both boundaries at once, and that the
available reporting cannot say whether the architecture on the other
side of either boundary behaves the way \cref{thm:selective-retention}
would recommend or simply degrades without anyone having designed for
the degradation.

\section{The Session-Trust Case}
\label{sec:session-trust}

One row resists easy classification, and it is worth working through
rather than forcing into the table, because it exposes a real gap in
how \cref{def:admissibility-relevant} was stated. Documented behavior
indicates that permission and trust decisions are not restored across
session boundaries: trust is re-established each session rather than
inherited from the last one. Read naively, this looks like
under-retention of an admissibility-relevant fiber --- whether a
directory or repository has been trusted plainly affects the
admissibility of later tool executions --- and \cref{thm:selective-retention}
would seem to predict a failure.

The naive reading is wrong, and the reason it is wrong matters more
than the specific case. \cref{def:admissibility-relevant} defines
relevance relative to a continuation family $C$
(\cref{def:continuation-family}), and $C$ is not fixed by the bare
fact that trust affects later behavior; it is fixed by what the system
is actually supposed to treat as a permitted continuation. If the
governing policy is that trust must be re-established every session
--- that a continuation in which stale, unre-verified trust is carried
forward is not itself admissible, regardless of what the trust state
used to be --- then the fiber in question is not admissibility-relevant
at all under the correctly specified $C$: no permitted continuation
ever depends on distinguishing which prior session's trust decision
occurred, because no permitted continuation is allowed to use a prior
session's trust decision in the first place. Discarding the state is
then not a retention failure but the direct implementation of the
admissibility structure itself.

This is a useful corrective to \cref{ch:falsifiability}'s own
apparatus: \cref{def:external-relevance}'s counterfactual test is only
as good as the continuation family it is run against, and getting $C$
wrong --- treating \enquote{trust used to be granted} as automatically
part of a permitted continuation, when the actual policy denies this
--- would misclassify a deliberate safety boundary as an engineering
oversight. The audit procedure does not resolve this automatically; it
requires knowing, independently, what the system's designers intended
$C$ to permit. Where that is documented, as it is here, the audit can
proceed correctly. Where it is not, this case is a warning about how
easily the same observation can be read two ways.

\section{A Case That Cuts the Other Direction}
\label{sec:cve-case}

\cref{ch:falsifiability} predicted, via
\cref{def:undecidability-failure} and
\cref{def:default-resolution-test}, that under-retention at a
genuinely admissibility-relevant fiber should show up not as a crash
but as a silent default resolution, exploitable if the default happens
to be predictable. A documented pre-patch vulnerability in the trust
pipeline matches this prediction closely enough to be worth stating as
a confirming instance rather than only a hypothetical one.

The relevant fiber is whether a request was authorized by a trust
decision the user genuinely made, or occurred before the user's trust
prompt was ever shown --- triggered instead by configuration supplied
through a malicious repository. Under the documented vulnerability,
the system could be induced to act on network configuration before the
trust-verification step had executed, meaning the two members of this
fiber --- genuinely-authorized request and request smuggled in ahead
of authorization --- were, at the point of execution, indistinguishable
to the system. This is \cref{def:undecidability-failure} exactly: a
downstream action (making the request) proceeds by defaulting, in
effect, to treating the fiber as though it had only one member. The
default-resolution test (\cref{def:default-resolution-test}) is
precisely what independent security researchers ran, informally, when
they demonstrated the exploit: constructing the two histories, holding
the eventual action fixed, and showing the system's behavior did not
track which one had actually occurred. Unlike the checkpoint case, and
unlike the compaction chain so long as its transcript survives, this
is not evidence that the architecture's retention tracks
admissibility-relevance; it is evidence of exactly
the failure mode the theory says under-retention at a relevant fiber
should produce, at a point in the pipeline where, prior to the patch
documented alongside the leak's disclosure, it evidently did not.

\section{What Would Have Falsified the Prediction Here}
\label{sec:leak-audit-falsification}

\cref{ch:falsifiability} named two falsification conditions. Neither
is decisively triggered by the material available for this audit, and
it is worth saying plainly what would have triggered them, so the
absence of falsification here reads as a real (if weak) test rather
than an inability to fail.

Falsification by uncorrelated retention would have required finding
that documented retention effort --- logging, checkpointing,
non-destructive history --- shows no tendency to concentrate at
documented non-injective transitions relative to injective ones: for
instance, if the architecture retained exhaustive logs of read-only,
side-effect-free operations (a case with no fiber to collapse in the
first place) at the same density as it retains state around
destructive edits and compactions. The available reporting shows the
opposite pattern --- retention concentrated at compaction boundaries,
pre-mutation checkpoints, and permission-relevant events, sparse
elsewhere --- which is consistent with, though it does not prove, the
correlation \cref{thm:selective-retention} predicts.

Falsification by silent success would have required finding a
documented case where a genuinely admissibility-relevant fiber was
under-retained and yet downstream judgments tracked the true
predecessor anyway, indicating some other mechanism was quietly doing
the work the audited channel was credited with. No such case appears
in the available reporting, though \cref{sec:leak-audit-limits}'s
caveats mean this is better read as \enquote{not found} than as
\enquote{ruled out.}

\section{Limits of This Audit}
\label{sec:leak-audit-limits}

The caveats from \cref{sec:audit-caveat} are worth restating in their
specific form here rather than left general. This chapter's
classifications rest on public technical reporting produced by
external parties reverse-engineering a leaked, minified-then-
reconstructed codebase, not on direct inspection of Anthropic's source
or design documentation. Several rows of \cref{tab:leak-audit} are
marked at reduced confidence for exactly this reason, and the
background daemon row in particular should be read as illustrative
rather than established, given that its existence is itself disputed
in the reporting this chapter draws on.

More importantly, the audit as conducted here cannot distinguish a
system whose retention architecture was \emph{designed} around
admissibility-relevance from one whose retention architecture merely
happens, for unrelated reasons, to correlate with it --- for instance,
if compaction boundaries are retained mainly because they are a
natural byproduct of how the summarization step is implemented, and
their admissibility-relevance is a fortunate side effect rather than a
design target. Settling that question would require access this
chapter does not have: internal design rationale, or a genuine
replay-based audit of the kind \cref{def:external-relevance} describes
in its strong form. What this chapter establishes is narrower and
still worth having: that the distinction between reconstruction-cost
and irreversibility materialization sorts a real, independently
documented architecture into two populated and behaviorally distinct
categories, that the sorting lines up with where the architecture is
reported to concentrate retention effort, and that at least one
documented failure --- \cref{sec:cve-case} --- matches the specific
failure signature the theory predicts for under-retention at a
relevant fiber, rather than some other signature the theory did not
anticipate.
% ===== END chapters/ch08-reading-the-leak =====
% ===== BEGIN chapters/ch09-transformer-semantics =====
\chapter{What the Leak Does Not Settle About Transformer Semantics}
\label{ch:transformer-semantics}

\cref{ch:category-mistake} left open the question that motivated this
entire monograph: whether a transformer's internal computation contains
anything functionally comparable to Logical Form, or whether its
apparently LF-like behavior --- tracking scope, binding, coreference,
and the rest --- is produced by some other means entirely. Having built
the reconstruction-cost and irreversibility apparatus, applied it to
HYDRA's own projection machinery, and audited a real leaked architecture
against it, we are in a position to say something sharper about this
question than \cref{ch:category-mistake} could: not an answer, but a
precise specification of what would count as one, together with an
important distinction the original debate ran together without
noticing.

\section{Two Bars, Not One}
\label{sec:two-bars}

The Logical Form debate, as inherited from the original public
argument and its rebuttals, implicitly asks a single question: does
the transformer \emph{have} an LF-like object. \cref{ch:hydra-projections}'s
apparatus makes visible that this is actually two questions, and that
conflating them is most of what has kept the debate unresolved.

\begin{definition}[Functional adequacy]
\label{def:functional-adequacy}
A system is functionally adequate with respect to an
admissibility-relevant fiber (\cref{def:admissibility-relevant}) if its
downstream behavior passes the audit procedure of \cref{ch:falsifiability}:
its outputs track the true predecessor at better than chance under the
default-resolution test (\cref{def:default-resolution-test}), whatever
internal mechanism produces that tracking.
\end{definition}

\begin{definition}[Architectural transparency]
\label{def:architectural-transparency}
A system is architecturally transparent with respect to a fiber if,
beyond being functionally adequate, it maintains an explicit,
inspectable object --- in the sense of \cref{def:estimate-as-set}, a
hypothesis set that can be read off directly rather than inferred by
probing --- whose state can be compared against the fiber without
requiring intervention on the system's internals.
\end{definition}

HYDRA's admissibility estimate, as formalized in
\cref{ch:hydra-projections}, is a design that targets both bars at
once: the hypothesis-set construction is meant to be functionally
adequate (it tracks $\pi^*$ under the conditions of
\cref{thm:convergence-characterization}) and, by construction,
architecturally transparent (the audit in \cref{ch:falsifiability} can
be run against $H_t$ directly). Nothing in the original Chomsky-derived
debate distinguished these two properties, and the transformer case is
exactly where they can come apart. A transformer could in principle be
functionally adequate over a given fiber --- its behavior genuinely
tracks the distinction, passing every version of the audit procedure
--- while remaining architecturally opaque, with no component anyone
has identified that plays the role $H_t$ plays for HYDRA. Conversely, a
system could be architecturally transparent by design and still fail
functional adequacy if its explicit machinery is wrong. The two are
independent, and settling one does not settle the other.

This reframes \cref{ch:category-mistake}'s open question. \enquote{Does
the transformer have an LF-analog} was really asking about
architectural transparency while gesturing at evidence --- correct
handling of scope, coreference, and the rest --- that only bears on
functional adequacy. A transformer could pass every behavioral test
that motivated the original question while the question, read
literally, remains open, because behavioral success is evidence for
\cref{def:functional-adequacy} and silent on
\cref{def:architectural-transparency}.

\section{A Transformer Audit Protocol}
\label{sec:transformer-protocol}

The audit procedure of \cref{ch:falsifiability} does not require
architectural transparency to be run; it requires only the ability to
construct controlled counterfactual histories and observe downstream
behavior, which is available for a transformer through ordinary prompt
intervention, without needing to interpret its internals at all.

\begin{enumerate}
\item \textbf{Identify a candidate fiber.} Choose a syntactic
construction whose surface form is compatible with more than one
underlying derivation --- exactly the cases
\cref{ch:category-mistake} used to motivate Logical Form in the first
place.
\item \textbf{Construct divergent predecessors.} Build two contexts
that supply different disambiguating evidence for which derivation is
intended, while keeping the ambiguous string itself, and everything
asked about it downstream, identical.
\item \textbf{Hold the probe fixed.} Ask the same downstream question
--- a paraphrase, an entailment judgment, a translation requiring the
distinction to be resolved --- against both contexts.
\item \textbf{Apply the default-resolution test.} Across many such
pairs, check whether the answers track the supplied disambiguating
evidence at better than chance
(\cref{def:default-resolution-test}). If they do, the transformer is
functionally adequate over that fiber. If they do not --- if the
answer is the same regardless of which context was supplied, or
tracks some surface feature uncorrelated with the intended derivation
--- it is not.
\item \textbf{Only then, and only if adequacy holds, ask about
transparency.} Attempt to localize a component of the computation ---
a direction in the residual stream, a subset of attention heads, an
intervention point discoverable by activation patching --- whose
value is injective over the fiber in the sense of
\cref{def:distinguishability}, independent of the final output. Success
here is evidence for architectural transparency; failure to find such
a component is not evidence against it, only evidence that it has not
been found, which is a weaker claim and should be reported as one.
\end{enumerate}

This protocol is offered as a methodological proposal, not as a result
this monograph has carried out. Its content, in the classification
this monograph has tried to be honest about throughout, is a framework
definition organizing future empirical work, not an established
finding.

\section{Mapping the Candidate Hypotheses onto Audit Outcomes}

Prior discussion of this question distinguished several ways a
transformer's LF-like behavior might arise. Restated against
\cref{def:functional-adequacy} and \cref{def:architectural-transparency},
each corresponds to a distinct, in-principle-distinguishable audit
outcome.

\begin{itemize}
\item \textbf{Genuine internal construction.} The transformer builds
something functionally isomorphic to an LF representation and
consults it. Predicted outcome: functional adequacy holds robustly,
including under adversarial variants of step 2 designed to break
surface correlations; and step 5 succeeds in localizing an injective
component, since a genuinely constructed intermediate representation
should leave some causal trace an intervention can find.
\item \textbf{Distributed approximation without an explicit object.}
The transformer's behavior tracks the distinction, but no
localized component carries it; the relevant information is smeared
across the representation in a way no single intervention point
isolates. Predicted outcome: functional adequacy holds robustly, as
in the first case, but step 5 fails to localize a component even under
sustained search --- functional adequacy without architectural
transparency, exactly the combination \cref{sec:two-bars} said the
original debate had no vocabulary for.
\item \textbf{Statistical shortcut mimicking the distinction.} The
transformer's behavior tracks the distinction only insofar as the
distinction correlates, in ordinary usage, with some surface feature
the model has learned to exploit --- word order frequency, lexical
co-occurrence, or similar. Predicted outcome: functional adequacy
holds under step 4's default form, but fails specifically under the
adversarial variant of step 2 that preserves the intended derivation
while breaking the surface correlation the shortcut depends on. This
is the outcome that would most directly vindicate the original
skeptical argument \cref{ch:category-mistake} was responding to, not
because a parser was found, but because behavior would fail exactly
where genuine derivation-sensitivity would be required and shortcut-
sensitivity would not suffice.
\end{itemize}

The three outcomes are empirically distinguishable from one another
using the same protocol, which is the protocol's main virtue: it does
not require deciding among them in advance, only running steps 1
through 5 and reading off which pattern of successes and failures
resulted.

\section{Worked Candidates}

\cref{ch:category-mistake}'s own examples are ready-made fibers for
step 1. \emph{I almost had a book stolen} supplies three derivations
distinguishable by which participant is agent, patient, or
would-be-agent of the theft; contexts can be constructed that
disambiguate toward each in turn, with a downstream probe (\enquote{who
was going to do the stealing}) that requires the distinction to be
resolved rather than merely gestured at. \emph{Every programmer fixed
a bug} supplies a scope ambiguity; contexts can be constructed
establishing either a single shared bug or one bug per programmer,
with a downstream probe (\enquote{how many distinct bugs were fixed})
that has a definite, derivation-dependent answer. Both are exactly the
kind of case \cref{def:external-relevance} needs: an unambiguous
downstream verdict whose correctness depends on which underlying
structure was intended, constructible without needing access to
anything beyond ordinary prompting.

\section{What the Leak Does Not Bear On}

It is worth being explicit that \cref{ch:reading-the-leak}'s audit and
this chapter's protocol concern entirely different objects. The leak
exposed the scaffolding around the model --- parsers, permission
systems, compaction machinery, checkpointing --- and
\cref{ch:reading-the-leak} audited that scaffolding's retention pattern
against \cref{thm:selective-retention}. Nothing in that audit touches
the transformer's own internal computation, which was never part of
what leaked; the model weights and forward-pass computation are not
source code, and no npm sourcemap exposes them. A reader who came away
from \cref{ch:reading-the-leak} thinking the leak had somehow settled
this chapter's question has conflated the two objects this section
is explicitly separating. The leak is evidence about architectural
decomposition in the surrounding system, exactly as
\cref{ch:category-mistake} argued from the outset; it was never
evidence, one way or the other, about what happens inside the model
itself.

\section{What This Chapter Establishes}

Not an answer. What it establishes is that the question motivating the
whole debate --- filtered through \cref{def:functional-adequacy} and
\cref{def:architectural-transparency} --- decomposes into two
independent, separately testable questions where the original
discussion had only one blurred question and no protocol for testing
it. It further establishes, via \cref{sec:transformer-protocol}'s
mapping, that the candidate hypotheses debated informally in the
original exchange are not merely different intuitions about the same
unresolvable mystery; they make different predictions about a concrete
and executable procedure, and the procedure does not require settling
questions about consciousness, understanding, or genuine semantics
that the informal debate kept circling back to. It requires only
running an intervention and reading off whether behavior tracks
derivation or tracks shortcut, and then, separately, whether anyone
can point to where.
% ===== END chapters/ch09-transformer-semantics =====

\part{Synthesis}
% ===== BEGIN chapters/ch10-hamiltonian-lagrangian-continuation =====
\chapter{Hamiltonian, Lagrangian, and Continuation Descriptions}
\label{ch:hamiltonian-lagrangian-continuation}

The apparatus built across Part II and Part III --- transition systems,
fiber collapse, reconstructive irreversibility, admissibility as a
relational property --- was developed without reference to physics.
This chapter places it against the two classical frameworks for
describing how systems move through state spaces at all, not as
imported vocabulary but as a way of saying precisely what this
monograph's framework is and is not doing. The point of the comparison
is not that continuation geometry resembles Hamiltonian or Lagrangian
mechanics; it is that the three frameworks answer three different
questions, and the difference is exact enough to state as propositions
rather than analogies.

\section{Hamiltonian Systems and State Sufficiency}

A Hamiltonian system evolves a state $(q, p) \in M$ --- position and
momentum coordinates on a phase space $M$ --- under Hamilton's
equations, generating for each $t$ a flow map $\Phi_t : M \to M$.

\begin{proposition}[Flow maps are diffeomorphisms]
\label{prop:hamiltonian-bijective}
For a Hamiltonian vector field smooth enough to guarantee existence
and uniqueness of solutions, the flow map $\Phi_t$ is, for each $t$
in its domain of definition, a diffeomorphism of $M$: smooth,
bijective, with smooth inverse $\Phi_{-t}$ (standard ODE existence
and uniqueness theory; see Arnold, 1989).
\end{proposition}

This is the entire content needed to place Hamiltonian systems within
the framework of \cref{ch:two-kinds-of-persistence}.

\begin{corollary}[Hamiltonian dynamics never collapse fibers]
\label{cor:hamiltonian-no-collapse}
Treating $(M, \emptyset, \Phi_t)$ as a transition system in the sense
of \cref{def:transition-system} (with a trivial input space, since
Hamiltonian evolution is autonomous), every fiber
$\mathrm{Fib}_P(y)$ under $\Phi_t$ is a singleton for every $y \in M$
and every admissible predecessor set $P$. Reconstructive
irreversibility (\cref{def:irreversibility}) therefore never occurs
in a Hamiltonian system: $\Phi_{-t}(y)$ is always a generating input
(\cref{def:generating-input}) for the predecessor state.
\end{corollary}

\begin{proof}
Immediate from bijectivity (\cref{prop:hamiltonian-bijective}): a
bijective map has fibers of size exactly one, so
\cref{prop:fiber-collapse}'s hypothesis is never satisfied, and
$\Phi_{-t}$, applied to $y$, recovers the unique predecessor exactly,
witnessing \cref{def:generating-input} directly.
\end{proof}

This is the precise sense in which a Hamiltonian system's present
state is sufficient for its past: not merely that reconstruction is
possible in some loose sense, but that \cref{def:irreversibility-materialization}
is vacuous for such a system by construction, and no history channel
is ever required on admissibility grounds
(\cref{thm:selective-retention}), because there is never a fiber for
one to be injective over. Liouville's theorem --- that Hamiltonian flow
preserves phase-space volume (Arnold, 1989) --- is the
measure-theoretic strengthening of this fact: not only is no
information lost pointwise, the flow does not even compress or expand
the space of possibilities locally.

\begin{remark}[Cost and information are still independent]
\label{rem:chaos-cost}
\cref{cor:hamiltonian-no-collapse} shows a Hamiltonian system is never
reconstructively irreversible; it says nothing about how expensive
reconstruction is. A chaotic Hamiltonian system --- one with positive
Lyapunov exponents --- can require reconstruction precision growing
exponentially in $t$ to recover $\Phi_{-t}(y)$ to any fixed tolerance,
making reconstruction-cost materialization
(\cref{def:reconstruction-cost}) of an approximate trajectory entirely
rational even though \cref{def:irreversibility-materialization} never
applies. This is \cref{prop:non-reducibility} given a physically real
instance: reconstruction can be made arbitrarily expensive without the
system ever crossing into the regime \cref{def:irreversibility}
describes, because expense and the existence of a generating input are
independent properties, exactly as \cref{prop:non-reducibility}
claimed in the abstract.
\end{remark}

\section{Lagrangian Systems and Trajectory Selection}

The Lagrangian formulation starts from a different primitive: not an
instantaneous state, but an entire trajectory $\gamma$, evaluated
through an action functional $S[\gamma] = \int L(\gamma(t),
\dot\gamma(t))\, dt$, with physical motion characterized by the
extremal condition $\delta S = 0$, yielding the Euler--Lagrange
equations (Arnold, 1989). Where the Hamiltonian question was
\emph{what can be recovered from a state}, the Lagrangian question is
\emph{which trajectory, among all those satisfying given boundary
conditions, is realized}.

\begin{remark}[Selection is orthogonal to reversibility]
\label{rem:lagrangian-orthogonal}
For standard classical Lagrangians, the Legendre transform relates the
Lagrangian formulation to a Hamiltonian one, and the resulting
dynamics inherit \cref{prop:hamiltonian-bijective}: the trajectory
selected by extremizing $S$ is itself a reversible flow in the sense
of \cref{cor:hamiltonian-no-collapse}. Extremal selection is therefore
not, by itself, a claim about reconstructive irreversibility; it is an
additional structure --- a global ranking over candidate trajectories
--- layered on top of dynamics that may or may not collapse fibers.
The Lagrangian picture's distinctive contribution is the idea of
\emph{selection by optimization}: a trajectory is picked out because
it extremizes a functional defined over the whole path, not because of
anything local to how each state transitions to the next.
\end{remark}

\cref{rem:lagrangian-orthogonal} isolates exactly the feature this
monograph's framework needs to contrast itself against: not
reversibility, which Part II already addressed directly, but
optimization as an explanatory principle. Admissibility
(\cref{def:admissibility}) has so far been defined without reference
to any functional to be extremized, and the next section makes that
omission a deliberate claim rather than an oversight.

\section{Continuation Systems and Admissible History}

\cref{def:admissibility} defines a state as admissible relative to a
trajectory and a continuation family --- a binary, relational
condition, not a ranking. This section states, and defends against an
obvious objection, the claim that this is a substantive difference
from the Lagrangian picture, not a restatement of it in different
notation.

\begin{proposition}[Admissibility does not require a global ranking]
\label{prop:no-ranking-required}
The admissibility relation $C$ of \cref{def:continuation-family} can
be specified, and \cref{thm:selective-retention} and
\cref{thm:historical-distinction} proved from it, without positing any
functional $S$ over trajectories such that $C(x)$ is characterized as
the extremizers of $S$. It suffices that $C(x)$ partition
continuations into permitted and forbidden, with no requirement that
permitted continuations be comparable to one another by any common
measure.
\end{proposition}

\begin{proof}[Justification]
Every definition and theorem from \cref{def:admissibility} through
\cref{thm:selective-retention} refers to $C(x)$ only through
set membership --- whether a given continuation lies in $C(x)$ ---
and never through a comparison between two members of $C(x)$. No step
in any proof in \cref{ch:formalizing-irreversibility} required
ranking, ordering, or measuring elements of $C(x)$ against each other.
\end{proof}

\begin{remark}[Answering the trivial-optimization objection]
\label{rem:trivial-optimization}
One might object that any binary partition can be trivially recast as
optimization of an indicator functional, $S[\gamma] = 0$ if $\gamma \in
C(x)$ and $S[\gamma] = -\infty$ otherwise, making
\cref{prop:no-ranking-required} empty. The objection is correct as far
as it goes and does not rescue the Lagrangian analogy: the substantive
content of the Lagrangian picture is not merely that \emph{some}
functional exists whose extremizers are the realized trajectories ---
under the indicator construction this is true of any deterministic or
admissibility-constrained system whatsoever --- but that the
functional is well-behaved enough to support the machinery
classical mechanics builds on it: smoothness, the Euler--Lagrange
equations, comparison of nearby trajectories by their action. The
indicator functional supports none of this; it is a relabeling of the
partition, not a variational structure. \cref{prop:no-ranking-required}'s
claim is that continuation geometry needs no such structure, trivial
or otherwise, to state or prove its central results, and the theorems
of \cref{ch:formalizing-irreversibility} are the evidence for this,
having been proved without it.
\end{remark}

With optimization set aside, what remains to characterize continuation
systems is exactly the machinery already built: transitions that
generically do collapse fibers (\cref{prop:repair-collapse}, in
contrast to \cref{cor:hamiltonian-no-collapse}), and persistence
secured not by selecting an optimal trajectory but by repair
--- revising an estimate under anomaly (\cref{def:repair}) so as to
remain within the admissible region, without any claim that the
resulting trajectory extremizes anything over its history.

\section{Three Regimes}

\begin{table}[htbp]
\centering
\footnotesize
\setlength{\tabcolsep}{4pt}
\begin{tabular}{p{2.2cm}p{3.0cm}p{3.0cm}p{3.0cm}}
\toprule
& \textbf{Hamiltonian} & \textbf{Lagrangian} & \textbf{Continuation} \\
\midrule
Primitive & instantaneous state $(q,p)$ & whole trajectory $\gamma$ & admissibility relative to a continuation family \\
Central question & what can be recovered from a state & which trajectory is realized & what survives reconstruction \\
Fiber collapse & never (\cref{cor:hamiltonian-no-collapse}) & inherited from underlying Hamiltonian dynamics; generically none & generic (\cref{prop:repair-collapse}) \\
Selection principle & none needed --- flow is deterministic and invertible & global extremization, $\delta S = 0$ & none required (\cref{prop:no-ranking-required}); persistence via repair under anomaly \\
\bottomrule
\end{tabular}
\caption{Three descriptions of trajectories through a state space, distinguished by what they take as primitive and what they require to be well-defined.}
\label{tab:three-regimes}
\end{table}

The three questions in \cref{tab:three-regimes}'s second row are the
chapter's central claim, restated once more because it is worth
having in a single place: Hamiltonian mechanics asks what can be
predicted from a state; Lagrangian mechanics asks which trajectory is
optimal; continuation geometry asks what survives reconstruction. The
first two questions are well-posed independently of anything in this
monograph and have been for centuries. The third is the question
\cref{ch:two-kinds-of-persistence} through \cref{ch:transformer-semantics}
have been answering, and \cref{tab:three-regimes} is intended to make
clear that it is not a special case, restriction, or generalization of
the other two, but a question with its own primitive --- admissibility
under a continuation family --- that is well-defined without borrowing
invertibility from the Hamiltonian picture or extremality from the
Lagrangian one.

\section{A Note on Variational Accounts of Persistence}

HYDRA's own account of its place in the wider literature notes an
approximate correspondence between the Free Energy Principle and
minimizing a combination of the RSVP accessibility potential and
entropy field.\footnote{The correspondence is stated at the level of
general resemblance, not formal equivalence, and is not re-derived
here; see the discussion of related work in HYDRA's own
presentation.} That correspondence, whatever its precise form, places
the Free Energy Principle on the Lagrangian side of
\cref{tab:three-regimes}: minimizing a variational quantity is
selection by optimization, in exactly the sense
\cref{rem:lagrangian-orthogonal} isolated. Continuation geometry's
departure from the Lagrangian picture in this chapter is therefore
also, to whatever extent the correspondence holds, a departure from
variational-optimization accounts of persistence generally, including
that one. This is noted here as a position within the wider landscape
of theories this framework sits alongside, not as a critique; nothing
in this chapter establishes that admissibility-based persistence is
correct where a variational account would be wrong, only that the two
are answering different questions, and that this monograph's
technical results (\cref{ch:formalizing-irreversibility} onward) did
not need the variational one to be stated or proved.
% ===== END chapters/ch10-hamiltonian-lagrangian-continuation =====
% ===== BEGIN chapters/ch10-continuation-as-resource =====
\chapter{Continuation as a Computational Resource}
\label{ch:continuation-as-resource}

The chapters of Part II and Part III were organized around a
distinction --- reconstruction-cost versus reconstructive-irreversibility
materialization --- applied first abstractly, then to HYDRA's own
architecture, then to a real leaked system, then to the question of
transformer semantics. This chapter asks what that distinction adds
up to when stated at its most general: not a classification tool for
sorting structures inside a particular architecture, but a claim about
a resource that computation depends on and that is not reducible to
the resources computer science already has names for. The claim is
stated as a theorem, not asserted as a slogan, because
\cref{prop:non-reducibility} already contains everything the proof
needs; what this chapter adds is the framing that makes the theorem
worth having stated at all.

\section{Three Resources and a Fourth Candidate}

Time, space, and energy are the resources computational complexity
theory is built around, and they share a property this chapter will
argue continuation does not: they are, within limits, fungible with
one another. Time-space tradeoffs are a standard subject of complexity
theory; energy-time tradeoffs govern physical implementations of
computation. Whatever the exact exchange rate, more of one resource
can typically compensate for less of another across a wide range of
problems.

The question this section poses is whether a system's capacity to
answer future admissibility judgments (\cref{def:admissibility})
belongs to this same family --- purchasable, in sufficient quantity,
by throwing more time, more space, or more energy at the problem --- or
whether it is a different kind of thing entirely.

\section{Distinction Budgets}

\begin{definition}[Distinction budget]
\label{def:distinction-budget}
Given a system's trajectory up to time $t$, its distinction budget
$B_t$ is the set of collapsed fibers $\mathrm{Fib}_P(y)$, for $y$
occurring at or before $t$, over which predecessor identity remains
recoverable --- either because a generating input still exists
(\cref{def:generating-input}) or because the system's retained history
channel is injective over that fiber (\cref{def:distinguishability}).
\end{definition}

\begin{proposition}[Monotonic shrinkage]
\label{prop:monotonic-shrinkage}
Absent a positive retention action at the moment a fiber collapses,
$B_t$ can only lose members as $t$ increases, never regain them: once
a fiber's generating input has ceased to exist and no injective
history channel was recorded for it at the time of collapse, no later
state of the system restores that fiber to $B_t$.
\end{proposition}

\begin{proof}
A fiber $\mathrm{Fib}_P(y)$ leaves $B_t$ exactly when neither a
generating input nor an injective retained channel is available for it
(\cref{def:distinction-budget}). Both conditions concern facts fixed at
or before the moment of collapse: the existence of a generating input
(\cref{def:generating-input}) is a fact about what remains
independently available in the world, and the injectivity of a
retained channel (\cref{def:distinguishability}) is a fact about the
specific $G$ applied at the transition itself
(\cref{def:augmented-system}). Neither fact can be altered by anything
that happens afterward: a generating input that has ceased to exist
cannot be un-destroyed, and a channel that failed to be injective at
the moment of retention cannot become injective retroactively, since
its value $h_{t+1}$ was already fixed by $G(h_t, x_t, u_t)$ at that
moment and nothing later changes what that value was.
\end{proof}

\begin{remark}
\cref{prop:monotonic-shrinkage} has the flavor of the data processing
inequality from information theory, which states that no processing
of a signal can increase the mutual information it carries about its
source. The resemblance is structural rather than formal --- this
chapter's setting concerns discrete distinguishability over
admissible predecessors rather than mutual information over a
channel, and no claim is made that \cref{prop:monotonic-shrinkage} is
a special case of the classical result --- but the shared shape is
worth noting: both say that a one-directional process can only destroy
distinguishing information already present, never manufacture it
after the fact.
\end{remark}

\section{The Resource Independence Theorem}

\begin{theorem}[Resource independence]
\label{thm:resource-independence}
Let $\mathrm{Fib}_P(y)$ be a fiber subject to reconstructive
irreversibility (\cref{def:irreversibility}). Then for any candidate
reconstruction procedure $P'$ and any allocation of time, space, or
energy to $P'$, no such allocation enables $P'$ to recover the
identity of the predecessor in $\mathrm{Fib}_P(y)$.
\end{theorem}

\begin{proof}
By \cref{def:irreversibility}, reconstructive irreversibility over
$\mathrm{Fib}_P(y)$ holds exactly when no generating input
(\cref{def:generating-input}) exists for the predecessor's identity.
A reconstruction procedure, however much time, space, or energy it is
allocated, is still a function from some input to an output; enlarging
its resource bound changes which functions are computable within that
bound, but does not change whether an input exists to compute from.
Since, by hypothesis, no state or artifact independent of the
system's own execution history determines the predecessor's identity,
there is no input, of any kind, from which any procedure --- regardless
of the time, space, or energy available to it --- could recover that
identity. The claim does not depend on complexity-theoretic bounds at
all; it depends only on the absence of a domain element to compute
from, which no resource allocation supplies.
\end{proof}

\cref{thm:resource-independence} is, in the classification this
monograph has tried to maintain throughout, close to definitional
given \cref{def:generating-input} and \cref{def:irreversibility} ---
its proof does not require anything beyond noticing that a
resource-bounded computation is still a computation from an input, and
that no allocation of resources creates an input where none exists.
Its value is exactly the value \cref{ch:falsifiability} assigned to
\cref{thm:selective-retention}: not surprise, but precision. It
converts \cref{prop:non-reducibility}'s abstract claim, and
\cref{rem:chaos-cost}'s physical illustration of it, into a direct
statement about the standard resource triad: continuation is not
purchasable with time, space, or energy, in the specific and limited
sense that no quantity of any of the three restores a distinction
whose generating input is gone.

\section{Why This Is Not Merely \enquote{More Memory}}

The most natural objection to treating continuation as a resource in
its own right is that it collapses into space --- into memory, the
already-recognized resource for retaining information across a
computation. The objection deserves a direct answer, because the
answer is what makes continuation a different kind of thing rather
than a fourth item on the same list.

\begin{proposition}[Continuation is not a quantity of memory]
\label{prop:not-memory}
Possessing unbounded memory does not, by itself, guarantee that a
system's retained structure is injective over any given
admissibility-relevant fiber. Injectivity (\cref{def:distinguishability})
is a property of \emph{what} is stored --- of the map $G$ applied at
the moment of collapse --- not of \emph{how much} storage is
available to store it in.
\end{proposition}

\begin{proof}[Justification]
A system with arbitrarily large memory can still apply a retention
policy $G$ that discards the specific information distinguishing two
predecessors --- storing, for instance, only that an edit occurred
and not what the edit was --- regardless of how much unused storage
capacity remains. \cref{def:distinguishability} depends only on
whether $G$ separates the fiber, which is a structural fact about $G$
independent of the size of $H$.
\end{proof}

This is why \cref{cor:retention-not-adequate} was necessary as early
as \cref{ch:two-kinds-of-persistence}: it is the specific instance of
\cref{prop:not-memory} that shows up as soon as retention is examined
concretely. Memory is a scalar resource, tradeable against time and
energy in the usual ways complexity theory studies. Continuation, at
the level of a single fiber, is not scalar at all; it is a binary fact
--- injective or not --- fixed by the structure of $G$ relative to that
fiber, and no amount of the scalar resource changes it. Aggregated
across many fibers, the \emph{breadth} of what a system's retention
preserves is measurable as $|B_t|$ (\cref{def:distinction-budget}),
which does behave something like a quantity; but $|B_t|$ grows only by
making $G$ injective over more fibers, which is a design choice about
what is stored, not a purchase made with additional storage capacity
on its own.

\section{Historical Observability, Defined}

The term this monograph's title and abstract have used informally
throughout can now be given a definition that ties directly to the
audit apparatus of \cref{ch:falsifiability}.

\begin{definition}[Historical observability]
\label{def:historical-observability}
Given a class $Q$ of possible future admissibility judgments a system
might need to make, its historical observability with respect to $Q$
at time $t$ is the proportion of $Q$ whose judgments remain decidable
given the current distinction budget $B_t$: the fraction of queries in
$Q$ that depend only on fibers still present in $B_t$, rather than on
fibers that have permanently exited it.
\end{definition}

\cref{def:historical-observability} makes precise what the audit
procedure of \cref{ch:falsifiability} was already measuring in
practice: the generating-input and fiber-injectivity tests are, in
effect, a way of sampling $Q$ and checking membership in $B_t$ without
needing to enumerate $Q$ exhaustively. A system's historical
observability is not a fixed property of its hardware or its raw
storage capacity; by \cref{thm:resource-independence}, it cannot be
increased after the fact by any allocation of time, space, or energy
once a relevant fiber has left $B_t$. It can only be preserved, at the
moment of collapse, by the specific choice --- \cref{def:repair-pruning}'s
Contraction and Soundness conditions, applied with an injective $G$
at exactly the admissibility-relevant fibers \cref{thm:selective-retention}
identifies --- to keep it from leaving in the first place.

This is the sense in which continuation is not merely another resource
alongside time, memory, and energy, but the resource that determines
which of a system's future admissibility judgments remain answerable
at all, independent of how much of the other three the system is
given. A system may have unlimited time to think, unlimited space to
store its thoughts, and unlimited energy to compute with, and still be
unable to answer a question whose answer depended on a distinction
its own dynamics discarded before anyone thought to ask it. Continuation,
in the sense this monograph has developed the term, is what stands
between a system and that specific kind of failure, and
\cref{thm:resource-independence} is the precise sense in which nothing
else substitutes for it.
% ===== END chapters/ch10-continuation-as-resource =====
% ===== BEGIN chapters/ch11-general-theory =====
\chapter{Toward a General Theory of Admissible Externalization}
\label{ch:general-theory}

Every use of \cref{thm:selective-retention} so far has been
retrospective: given a system that already exists, audit its retention
pattern against admissibility-relevance and classify what is found.
This chapter turns the criterion around. If a designer is building a
system from scratch, under a finite budget for retained history, what
does the theory developed in Part II and Part III actually recommend,
and how far beyond HYDRA and the Claude Code leak does the
recommendation generalize?

\section{From Audit Criterion to Design Principle}

\cref{thm:selective-retention} states a necessary and sufficient
condition for a retention architecture to support exactly the
admissibility judgments a system's continuation families require, at
minimum cost. Read prospectively, this is a specification: a designer
who can enumerate the admissibility-relevant fibers
(\cref{def:admissibility-relevant}) their system's dynamics will
generate has, in principle, a target to build toward --- injective
history channels exactly there, and nowhere else.

The qualification \enquote{in principle} is doing real work. Real
systems do not have unlimited capacity to build injective channels
everywhere \cref{thm:selective-retention} would credit them for
building one. The theorem is silent on what to do when the set of
admissibility-relevant fibers exceeds what a bounded retention budget
can cover, and that silence is this chapter's starting point.

\section{The Allocation Problem}

\begin{definition}[Retention budget and candidate fibers]
\label{def:allocation-problem}
Let $F_1, \dots, F_n$ be candidate fibers a system's designer has
identified as plausibly admissibility-relevant, each with an estimated
relevance weight $w_i > 0$ (reflecting the severity of the
undecidability failure, \cref{def:undecidability-failure}, that would
result from under-retention at $F_i$) and a cost $c_i > 0$ of building
an injective history channel over $F_i$. Given a total retention
budget $\beta$, the allocation problem is to choose a subset $S
\subseteq \{1, \dots, n\}$ maximizing $\sum_{i \in S} w_i$ subject to
$\sum_{i \in S} c_i \leq \beta$.
\end{definition}

\begin{remark}[This is not a redefinition of admissibility]
\label{rem:not-redefining-admissibility}
\cref{def:allocation-problem} introduces weights and a budget
constraint, and it is worth being precise about what is and is not
being optimized. \cref{ch:hamiltonian-lagrangian-continuation}
established that the admissibility relation $C$ itself
(\cref{def:continuation-family}) requires no extremal structure to be
well-defined, and nothing here revises that: $C$ is fixed before this
section begins, and $F_1, \dots, F_n$ are fibers that
\cref{def:admissibility-relevant} --- itself defined without reference
to any weighting --- has already picked out as relevant relative to
that fixed $C$. What is being optimized is a downstream engineering
decision: given that these fibers matter and capacity is limited,
which to prioritize. This is a claim about resource allocation under a
fixed admissibility structure, not a claim that admissibility is
itself the output of an optimization, and
\cref{rem:trivial-optimization}'s distinction between a genuine
variational structure and mere resource bookkeeping applies again
here in the designer's favor: nothing about $C$ has been asked to be
smooth, comparable, or extremal.
\end{remark}

\cref{def:allocation-problem} is an instance of the knapsack problem,
well known to be NP-hard in general (standard result; see Cormen,
Leiserson, Rivest, and Stein, \emph{Introduction to Algorithms}, or
any comparable algorithms text treating weighted knapsack). This is
stated not to claim a new complexity result but to be honest about
what has and has not been achieved: the allocation problem has been
given a precise combinatorial shape, which it did not have before this
chapter, but a precise shape is not the same as a tractable one, and
\cref{sec:general-open} returns to this directly.

\section{A Design Principle: Admissibility-Aware Externalization}

Where exact solution of \cref{def:allocation-problem} is infeasible,
the structure of the problem still licenses a design heuristic, in the
same spirit as the greedy approximation ratio available for weighted
knapsack: prioritize candidate fibers by relevance-per-cost, $w_i /
c_i$, building injective channels for the highest-ratio candidates
first until the budget is exhausted.

\begin{itemize}
\item[] \textbf{Admissibility-Aware Externalization (design principle).}
Given finite retention capacity, a system should allocate history
channels to admissibility-relevant fibers in order of relevance per
unit retention cost, rather than uniformly, rather than in order of
raw reconstruction difficulty (\cref{def:reconstruction-cost}, which
\cref{prop:non-reducibility} already showed is the wrong axis), and
rather than by implementation convenience.
\end{itemize}

This is stated as a design principle, not a theorem, because its
force depends on the accuracy of the weights $w_i$ and costs $c_i$ a
designer supplies, which \cref{ch:falsifiability}'s apparatus can help
estimate --- $w_i$ via the severity implied by
\cref{def:undecidability-failure} at $F_i$, $c_i$ via the ordinary
engineering cost of an injective $G$ over $F_i$ --- but cannot
guarantee are correct in advance of deployment. \cref{cor:two-failure-modes}'s
two failure modes are exactly what a designer following this principle
imperfectly will produce: under-retention where $w_i$ was
underestimated, over-retention where it was overestimated, both
detectable after the fact by the audit procedure even when they were
not foreseeable before it.

\section{Generalization Beyond HYDRA}

Every application so far --- \cref{ch:hydra-projections}'s formal
apparatus, \cref{ch:reading-the-leak}'s audit,
\cref{ch:continuation-as-resource}'s resource-independence argument
--- has concerned computational systems narrowly construed: software
architectures with explicit transition dynamics. HYDRA's own
presentation makes a considerably broader claim, proposing that the
same quartet of trajectory space, projection system, admissibility
structure, and persistence functional recurs across cognition,
biology, institutions, physics, and cosmology, with the Structural
Universality Conjecture as its most general statement. This monograph
does not attempt that scope. What it can state, at a scope matched to
what \cref{ch:formalizing-irreversibility} through
\cref{ch:continuation-as-resource} actually established, is narrower
and specifically about retention.

\begin{itemize}
\item[] \textbf{Admissible Externalization Conjecture.} Any persistent
system operating under finite retention capacity, whose transition
dynamics are generically non-injective
(\cref{def:fiber}) and whose continuations are governed by an
admissibility structure in the sense of \cref{def:admissibility}, will
--- to the extent its retention architecture has been shaped by
selective pressure toward correct future functioning rather than by
accident --- exhibit a retention pattern correlating with
admissibility-relevance (\cref{def:admissibility-relevant}) rather
than with raw reconstruction cost, and departures from this
correlation will be observable, via the audit procedure of
\cref{ch:falsifiability}, as the specific failure signatures of
\cref{def:undecidability-failure} and over-retention.
\end{itemize}

Following the classification this monograph has used throughout: this
is a conjecture, not a theorem, and it is falsifiable in exactly the
sense \cref{ch:falsifiability} made precise, applied now across
domains rather than within one architecture. Three illustrative
candidates, offered as motivating instances rather than as evidence
for the conjecture, suggest where a genuine test might be conducted.
Legal and financial audit trails are retained selectively around
transactions and decisions with downstream liability consequences, not
uniformly across all record-generating activity, which is the pattern
the conjecture predicts if liability determination is treated as an
admissibility judgment. Version control systems retain commit history
essentially unconditionally rather than selectively, which --- if the
conjecture is to survive this case --- would need to be explained by
the extremely low marginal cost of that retention ($c_i \approx 0$ for
nearly all fibers) rather than by high $w_i$ everywhere, an alternative
the allocation problem (\cref{def:allocation-problem}) already
accommodates, since low-cost, low-relevance fibers can still be worth
including when $\beta$ is not binding. Biological memory consolidation
is reported in the relevant literature to favor emotionally salient or
consequential events over routine ones, which would instantiate the
conjecture if salience tracks something like admissibility-relevance
to future decisions, though establishing that correspondence
rigorously is well outside this monograph's scope and is flagged here
only as a direction, not a claim.

\section{What Remains Open}
\label{sec:general-open}

Four gaps are worth naming plainly, in keeping with this monograph's
practice of distinguishing what has been established from what has
only been proposed.

First, \cref{def:allocation-problem}'s NP-hardness means the design
principle of \cref{ch:general-theory} is a heuristic with the usual
gap between greedy approximation and exact optimality; nothing here
establishes how large that gap is for retention-allocation problems
specifically, as opposed to knapsack problems in general.

Second, the weights $w_i$ and costs $c_i$ the allocation problem
requires are, in any real system, estimates rather than known
quantities, and this chapter has not supplied a method for estimating
them beyond gesturing at \cref{ch:falsifiability}'s apparatus.

Third, \cref{sec:hydra-scope}'s open question --- whether a system's
hypothesis space $\mathcal{A}$ is fixed in advance or expands as
evidence arrives --- bears directly on the allocation problem too: a
designer enumerating candidate fibers $F_1, \dots, F_n$ in advance
implicitly assumes the space of possible fibers is known at design
time, which may fail for the same reasons \cref{sec:hydra-scope}
flagged for $\mathcal{A}$.

Fourth, and most directly, the Admissible Externalization Conjecture
has been illustrated, not tested. \cref{ch:reading-the-leak} conducted
one specification-level audit of one system; the conjecture's claim is
about a pattern across systems and domains that this monograph has not
attempted to survey. The conjecture is stated in a form intended to
be testable by exactly that kind of survey, using the audit procedure
already built, rather than in a form that would require new
theoretical machinery to evaluate --- but the survey itself remains
undone.
% ===== END chapters/ch11-general-theory =====
% ===== BEGIN chapters/ch11b-variational-primacy =====
\chapter{Variational Primacy and the Missing Axis}
\label{ch:variational-primacy}

\cref{ch:hamiltonian-lagrangian-continuation} contrasted continuation
geometry against Hamiltonian and Lagrangian mechanics as those
frameworks are ordinarily taught, with Lagrangian mechanics presented
as, in effect, an alternative derivation of Newton's laws. A natural
objection is that this undersells the Lagrangian picture: in its most
philosophically serious form, Lagrangian mechanics is not a derivation
of anything more fundamental, but the more fundamental framework in
its own right, with Newtonian mechanics as a restrictive special case.
This chapter takes that objection seriously against an actual text
that argues exactly this, checks whether the distinction
\cref{ch:hamiltonian-lagrangian-continuation} drew survives contact
with the strongest available version of the variational picture, and
finds a further, more precise point of contrast along the way that the
weaker version of the comparison would have missed.

\section{A Text That Already Resists the Naive Reading}

José Rachid Mohallem's \emph{Lagrangian and Hamiltonian Mechanics: A
Modern Approach with Core Principles and Underlying Topics}
(Springer, 2024) opens its preface with a single sentence set apart
from the surrounding text: \enquote{In the beginning, there was the
Action Principle.} The book's stated aim is to correct a specific
pedagogical habit --- presenting the Lagrangian of a conservative
system as $L = T - V$ from the outset and treating Hamilton's
Principle as, in the author's words, an alternative way of obtaining
the results of Newtonian Mechanics --- in favor of treating the
search for Lagrangians from symmetries, via Hamilton's Principle, as
what the author calls the more powerful theoretical instrument of
physics. This is not the textbook-standard framing
\cref{ch:hamiltonian-lagrangian-continuation} contrasted itself
against; it is a considerably stronger position, treating variational
selection as explanatorily prior to, rather than derivable alongside,
Newtonian dynamics.

\section{Confirming the Absence}

If any single text were likely to have already developed something
functionally close to this monograph's distinguishability axis, a
philosophically self-aware treatment organized explicitly around
Hamilton's Principle as primitive, rather than around $F = ma$, would
be a reasonable place to look. A search of the complete extracted text
of Mohallem's book --- all 152 pages --- for any of the terms
\emph{irreversib-}, \emph{reversib-}, \emph{predecessor}, or
\emph{distinguishab-} returns no occurrences whatsoever. This is
reported as an observation about one text, not a survey of the
literature on classical mechanics, and no claim is made here that no
treatment anywhere engages a comparable question. What can be said is
narrower and still useful: even a text that explicitly rejects the
reduction of Hamilton's Principle to a mere alternative route to
Newton's laws, and that spends a full chapter on the Hamiltonian
formulation's own conceptual apparatus, organizes its entire
discussion around two questions --- which trajectory is selected, and
what follows from a given state --- and at no point asks which
predecessor states remain distinguishable from a given successor.
\cref{tab:three-regimes}'s third column is not competing with a
question this text already asks under different vocabulary; the
question does not appear.

\begin{remark}[A closer near-miss than the bare search suggests]
\label{rem:determinism-near-miss}
A lexical search of this kind is blunter than it looks, and it is
worth naming the place it comes closest to failing. The book does
discuss \emph{determinism} directly, in a passage noting that
classical mechanics' determinism of motion is among the paradigms
replaced when passing to a quantum treatment. This is, in substance,
an informal statement of state-sufficiency --- the Hamiltonian side of
\cref{tab:three-regimes} --- under vocabulary the original four-term
search did not anticipate. It does not, on inspection, extend to this
monograph's question: the passage is entirely forward-looking, about
what a state determines about the future, and at no point turns the
question around to ask whether a state determines its past to the
same degree. The distinction survives, but the search that first
suggested it was cruder than it should have been, and a reader should
take the absence of the four searched terms as a starting point for
\cref{rem:determinism-near-miss}'s closer reading, not as a
substitute for it.
\end{remark}

\section{Conservation Is Not Distinguishability}

The comparison is worth pushing one step further than a text search,
because the book's own central technical apparatus --- constants of
motion, arising from symmetries of the Lagrangian via cyclic
coordinates --- initially looks like it might be a variant of this
monograph's distinguishability-preservation concept wearing different
notation. It is not, and the precise way it is not is worth stating
formally, because the two are not merely different but, in a specific
sense, structurally opposed.

\begin{definition}[Constant of motion]
\label{def:constant-of-motion}
A quantity $F(q, \dot q, t)$ is a constant of motion along a
trajectory if its total time derivative vanishes identically along
that trajectory: $\dot F(q, \dot q, t) = 0$. Constants of motion
typically arise, via Noether's theorem in its classical-mechanical
form, from a symmetry of the Lagrangian --- a cyclic coordinate not
appearing explicitly in $L$ yields a conserved conjugate momentum.
\end{definition}

\begin{proposition}[Constants of motion are generically non-injective on trajectory space]
\label{prop:conservation-noninjective}
Let $F$ be a constant of motion and let $\mathcal{T}$ be the space of
dynamically realizable trajectories. The level sets of $F$ --- the
fibers of the map sending each trajectory to its (constant) value of
$F$ --- generically contain more than one distinct trajectory: except
in fully integrable systems with as many independent constants of
motion as degrees of freedom, knowing the value of $F$ alone does not
determine which trajectory, among those sharing that value, is
realized.
\end{proposition}

\begin{proof}[Justification]
This is the ordinary content of partial integrability: a single
constant of motion, such as energy, restricts a trajectory to an
energy shell of the phase space, generally a submanifold of dimension
one less than the full phase space, and this restriction is
compatible with a continuum of distinct trajectories within that
shell whenever the system has more than one degree of freedom and
fewer independent constants of motion than are needed to fix a
trajectory uniquely.
\end{proof}

\cref{prop:conservation-noninjective} is the precise sense in which
tracking a constant of motion is the structural opposite of what
\cref{def:distinguishability} asks a history channel to do. A history
channel earns its keep, in this monograph's framework, by being
\emph{injective} over an admissibility-relevant fiber --- by
separating predecessors that would otherwise be conflated. A constant
of motion is, by construction, a coarsening: it is valuable to physics
precisely because it is the same for many different trajectories,
which is what makes it a useful invariant to reason with rather than a
complete description of the motion. Where this monograph's history
channels are built to discriminate, constants of motion are built to
forget everything except one shared feature.

\begin{remark}[Conservation tracking is reconstruction-cost compression]
\label{rem:conservation-is-cost-side}
\cref{cor:hamiltonian-no-collapse} already established that a
Hamiltonian system's raw state never loses distinguishing information
under its own dynamics: the flow map is bijective, so nothing is ever
irreversibly lost at the level of $(q,p)$ itself. An observer who
tracks only a constant of motion is therefore never doing anything
this monograph would classify as irreversibility materialization
(\cref{def:irreversibility-materialization}) --- the full state remains,
in principle, recoverable throughout, exactly as \cref{rem:chaos-cost}
observed can remain true even when reconstruction is computationally
severe. Tracking $F$ instead of the full state is a reconstruction-cost
choice (\cref{def:reconstruction-cost}): a deliberate, tractability-
motivated compression of a quantity that was never, in the relevant
sense, at risk of becoming unrecoverable, because within a genuinely
Hamiltonian regime nothing is. Classical conservation laws, read
through this monograph's classification, sit entirely on the
reconstruction-cost side of \cref{tab:three-regimes}, never on the
continuation side --- which is a further reason \cref{tab:three-regimes}'s
third column was not available to a text organized, however
sophisticatedly, around constants of motion and variational selection:
its central technical tool is, on this framework's own terms, an
instance of the axis \cref{prop:non-reducibility} already distinguished
from the one this monograph is about.
\end{remark}

\section{What This Adds to Chapter 10}

\cref{tab:three-regimes} stands as stated: the Hamiltonian column's
claim that no history channel is ever required needs no revision, and
neither does the Lagrangian column's claim that extremal selection is
orthogonal to reversibility. What this chapter adds is a stronger
warrant for the claim that continuation geometry's organizing question
is not merely absent from the \emph{textbook} presentation of
classical mechanics, but absent from at least one serious attempt to
present the subject with variational selection as fully primary rather
than as a derived convenience --- and that the reason it is absent is
not oversight but structural: the tool this stronger presentation
relies on most heavily, the constant of motion, answers a
coarsening question, and \cref{ch:two-kinds-of-persistence}'s
distinguishability question needed a different technical apparatus
because it is asking in the opposite direction.

\section{A Remaining Caveat}

One book, however carefully argued, is one data point. Mohallem's text
was consulted here because it was made available for exactly this
comparison and because its stated thesis is unusually well-suited to
testing whether \cref{ch:hamiltonian-lagrangian-continuation}'s
contrast survives the strongest available version of the objection.
Nothing in this chapter should be read as a claim that no treatment of
classical or analytical mechanics, anywhere, engages a question
resembling this monograph's distinguishability axis; only that this
one, chosen because it was likely to be the hardest case available,
does not.

A specific and more promising place to look, named but not checked
when this chapter was first drafted, turns out on investigation to
matter more than a passing caveat can convey, and the earlier version
of this section understated the case. Constrained Hamiltonian systems
--- those in which the Legendre transform from velocities to momenta
is singular, so that the map fails to be invertible --- are a
well-known site of genuinely non-injective structure inside classical
mechanics itself, formalized in the Dirac--Bergmann theory of
constraints. The kernel of the singular Legendre transform is, by
construction, exactly a fiber in the sense of \cref{def:fiber}: distinct
velocity configurations mapping to the same momentum. This is a
substantively different phenomenon from anything
\cref{cor:hamiltonian-no-collapse} addresses, since that result
concerned the time-evolution flow map for a fixed, non-degenerate
Hamiltonian, not the map between Lagrangian and Hamiltonian
descriptions at fixed time. Mohallem's book does not treat this case.
The wider literature does, at length, and the treatment is closer to
this monograph's question than a first pass suggested.

\section{The Problem of Time and the Problem of Observables}

The standard resolution of a singular Legendre transform's
non-injectivity is to declare it physically empty: points related by
the transformations generated by the resulting first-class constraints
are said to represent the \emph{same} physical state, foliating phase
space into \enquote{gauge orbits,} and a \emph{Dirac observable} is
defined as a phase-space function invariant along these orbits ---
formally, one whose Poisson bracket with every first-class constraint
vanishes. This is, in the vocabulary \cref{ch:hydra-projections}
already established, a projection: a map constant on the fibers
induced by gauge equivalence, in exactly the sense
\cref{def:projection} used the word before this chapter existed to
complicate it. Read this way, gauge theory's resolution of its own
non-injectivity is to stipulate, at the level of the theory itself,
that the collapsed distinctions were never physical to begin with ---
not to ask, as \cref{def:admissibility-relevant} does, whether some
downstream judgment depends on them.

This stipulation is not as settled as the standard presentation makes
it sound, and general relativity is where it visibly breaks. Because
the Hamiltonian of general relativity in its constrained form is
itself a sum of first-class constraints, the standard argument implies
the Hamiltonian generates gauge transformations rather than physical
evolution, so that Dirac observables --- gauge-invariant by
construction --- do not evolve with respect to it at all. This
produces the \emph{problem of time}: formally, nothing happens, while
observationally, quantities such as cosmological redshift plainly do
change. The literature's response has not converged. Some authors
maintain the orthodox reading, that the Hamiltonian constraint is pure
gauge and apparent evolution must be reconstructed some other way; others
argue the identification of the Hamiltonian constraint as a pure gauge
generator is mistaken --- one recent treatment states directly that
Dirac's own argument for this identification does not hold in general,
following an extended line of criticism running through Bergmann,
Kuchař, Pons, Salisbury, and Anderson. Whichever position is correct,
the dispute itself is the relevant fact for this chapter: physicists
working on exactly this problem are, in substance, asking whether a
phase-space distinction collapsed by a constraint carries physical
content or does not, and have not settled it after decades of direct
engagement.

\begin{remark}[What is actually different]
\label{rem:gauge-vs-admissibility}
Stated this plainly, the problem of observables in constrained
Hamiltonian systems is close enough to this monograph's question that
\cref{ch:variational-primacy}'s original claim --- that no treatment
of classical mechanics asks anything resembling it --- was too strong,
and is withdrawn in that form. What survives is a narrower and more
precise distinction between the two questions, not their identity.
The Dirac-observable criterion, $\{F, C_a\} = 0$ for every first-class
constraint $C_a$, is a fixed, algebraic, theory-level test: a
distinction either is or is not gauge, once and for all, independent
of any downstream use. \cref{def:admissibility-relevant}'s criterion
is explicitly relative to a continuation family $C$
(\cref{def:continuation-family}): the same fiber can be
admissibility-relevant under one $C$ and irrelevant under another, as
\cref{sec:session-trust} demonstrated directly for the cross-session
trust fiber. Physics is attempting to answer, for general relativity,
a single context-free version of the question this monograph's
apparatus only ever asks relative to a specified continuation family.
\end{remark}

\begin{remark}[A speculative connection, flagged as such]
\label{rem:problem-of-time-speculation}
It is tempting, and worth stating explicitly as speculation rather
than as anything this monograph establishes, that
\cref{rem:gauge-vs-admissibility}'s distinction may bear on why the
problem of time has resisted decades of direct effort by researchers
considerably more expert in the relevant physics than this monograph
can claim to be. If \cref{def:admissibility-relevant}-style relevance
is genuinely relative to a continuation family and has no
context-independent fact of the matter in general, then a research
program seeking a single, theory-level, context-free answer to
\enquote{does this constraint-generated distinction matter} would be
expected to encounter exactly the kind of persistent, foundational
disagreement the literature shows, not because the physics is
unusually difficult in the ordinary sense, but because the question,
posed context-free, may not have the kind of answer being sought. This
is offered only as a possible reframing, not a resolution; nothing in
this monograph's apparatus is equipped to adjudicate a live dispute in
the foundations of canonical general relativity, and no such
adjudication is attempted here.
\end{remark}

This qualifies \cref{ch:hamiltonian-lagrangian-continuation} and the
earlier sections of this chapter rather than undermining them. The
axis \cref{tab:three-regimes} identifies is not simply absent from
physics; it surfaces exactly where constrained Hamiltonian systems
push hardest against the gauge-orbit resolution, in one of theoretical
physics's most famously unresolved foundational problems. That the
question turns out to be live rather than absent is, if anything, a
stronger indication that it is a real question than the original,
overstated version of this chapter claimed.
% ===== END chapters/ch11b-variational-primacy =====
% ===== BEGIN chapters/ch12-conclusion =====
\chapter{Conclusion}
\label{ch:conclusion}

The argument began with a narrow question: whether finding a parser
inside Claude Code's leaked source told us anything about whether the
model surrounding it understands code. \cref{ch:category-mistake}
answered that question in one chapter and did not need the apparatus
built afterward to do it. Everything from \cref{ch:why-structure-externalizes}
onward exists because a more interesting question was sitting
underneath the one the public debate was asking: not whether systems
decompose into learned and non-learned parts, which was never
seriously in doubt, but whether there is a principled account of which
computations get externalized as retained structure and which do not.
This chapter states what was actually established in pursuit of that
question, classifies each claim by how much weight it can bear, and
says plainly where the argument could turn out to be wrong.

\section{What This Monograph Established}

Four pieces of apparatus, built in sequence, carry the argument.

\cref{ch:two-kinds-of-persistence} separated two reasons a structure
persists --- reconstruction cost and reconstructive irreversibility ---
and showed, via the Historical Distinction Theorem
(\cref{thm:historical-distinction}), that the two are not points on a
shared scale: one concerns whether a generating input exists, the
other concerns whether a retained channel is injective over a
collapsed fiber, and no amount of the first substitutes for the
second.

\cref{ch:formalizing-irreversibility} made admissibility relational,
in the sense HYDRA's own architecture requires, and proved the
Selective Retention Theorem (\cref{thm:selective-retention}): a system
need only preserve distinguishability where some future admissibility
judgment actually depends on it, not everywhere reconstruction happens
to be impossible. \cref{ch:hydra-projections} then gave the estimate
side of this apparatus a precise treatment as hypothesis pruning,
proved exactly when repair converges to the truth rather than merely
stabilizing (\cref{thm:convergence-characterization}), and exhibited a
constructed case where it does not.

\cref{ch:falsifiability} turned the theorem into something checkable:
an independence criterion for admissibility-relevance
(\cref{def:external-relevance}) that does not consult a system's
retention architecture to determine what that architecture should
contain, an audit procedure built on it, and two named conditions
under which the whole apparatus would be considered wrong for a given
system rather than merely inconvenient.

\cref{ch:continuation-as-resource} gave the argument its most general
statement: the Resource Independence Theorem
(\cref{thm:resource-independence}), showing that no allocation of
time, space, or energy substitutes for an injective history channel
once its generating input is gone, which is the precise sense in which
continuation is not a fourth member of the usual resource family but a
different kind of thing.

Everything else --- the leak audit, the HYDRA comparison, the
transformer-semantics reframing, the classical-mechanics contrast, the
design principle and its generalization --- is these four results
doing work against specific material: a real leaked architecture, a
formal cognitive architecture, an open question in AI, and two
carefully chosen physics texts.

\section{The Status of These Claims}

Not everything in this monograph carries the same evidentiary weight,
and the classification below is offered so that no claim is
mistakenly read as stronger than what was actually shown.

\begin{itemize}
\item \textbf{Proved from stated definitions, close to tautological
given them.} The Historical Distinction Theorem, the Selective
Retention Theorem, the Non-Reducibility and Resource Independence
propositions, and the Convergence Characterization Theorem. Their
proofs are short because their content is mostly already present in
the definitions; their value, as stated repeatedly throughout, is
organizational rather than surprising. None of these can be falsified
by any observation, because none of them makes a claim about the
world independent of the definitions that generate it.
\item \textbf{Framework definitions, not claims.} Admissibility,
continuation families, projections, generating inputs, fibers,
distinguishability preservation, historical observability. These are
not true or false; they are the vocabulary the rest of the monograph's
claims are stated in, and their only defense is that they turned out
to be usable --- that real distinctions (compaction, checkpoints, trust
state) could be sorted with them without the sorting collapsing into
vacuity.
\item \textbf{An empirical claim about one real system, checked at
reduced confidence.} \cref{ch:reading-the-leak}'s audit of the Claude
Code leak. This is the closest thing in the monograph to a test
against reality, and it was explicit throughout about being a
specification-level audit --- sampled, not exhaustive; drawing on
public reporting, not source access; producing at least one case
(context compaction) the audit could not cleanly resolve rather than
forcing a resolution.
\item \textbf{A proposed protocol, not a result.} \cref{ch:transformer-semantics}'s
audit procedure for transformer semantics. Nothing in this monograph
ran it. Its contribution is decomposing one blurred question into two
checkable ones, not answering either.
\item \textbf{A conjecture, explicitly falsifiable, not tested beyond
illustration.} The Admissible Externalization Conjecture
(\cref{ch:general-theory}). The candidate domains offered alongside it
--- audit trails, version control, memory consolidation --- were named
as motivating instances, not evidence.
\item \textbf{A comparison against specific, named sources, not a
literature survey.} \cref{ch:hamiltonian-lagrangian-continuation} and
\cref{ch:variational-primacy}'s contrast with classical mechanics,
checked against one textbook chosen for being the hardest available
case, plus a follow-up check of the constrained-Hamiltonian-systems
literature that partially retracted the chapter's original claim: the
distinguishability question turns out to be live, under the vocabulary
of gauge orbits and Dirac observables, in the unresolved problem of
time in canonical general relativity, though the specific
context-relative reframing \cref{rem:problem-of-time-speculation}
offers remains explicitly speculative. The determinism passage named
in \cref{rem:determinism-near-miss} remains a near miss rather than an
instance of the question.
\end{itemize}

\section{A Note on Terminology}
\label{sec:terminology-migration}

One further piece of accounting belongs here rather than in any single
chapter, because it concerns the whole document's vocabulary rather
than any specific claim. \emph{Reconstructive irreversibility} and
\emph{irreversibility materialization} are the names
\cref{ch:two-kinds-of-persistence} gave the phenomenon this monograph
studies, and they are the names that connect the technical apparatus
back to the motivating question in the title. They are not, however,
the vocabulary the load-bearing proofs are actually stated in.
\cref{ch:formalizing-irreversibility}'s Selective Retention Theorem
and \cref{ch:hydra-projections}'s convergence results --- the two
chapters doing the most proof-theoretic work in the book --- are
stated and proved entirely in terms of fibers, injectivity, and
distinguishability preservation, and do not cite the irreversibility
definitions directly at any point; neither does
\cref{ch:transformer-semantics}'s protocol. The irreversibility
vocabulary is concentrated instead in \cref{ch:two-kinds-of-persistence}
itself, where it is defined; in \cref{ch:reading-the-leak}, where it
functions as a classification label for the audit table; and in
\cref{ch:hamiltonian-lagrangian-continuation} and
\cref{ch:variational-primacy}, where it is the named category being
contrasted against classical mechanics.

This is not a defect, and no revision follows from it: renaming the
central term to track the apparatus precisely --- something like
\enquote{distinguishability materialization} --- would be more
technically accurate and would cost the document its connection to the
question that motivated it in the first place, since \enquote{what
survives reconstruction} is a question about irreversibility, asked in
the vocabulary a reader arrives with, and the technical machinery of
fibers and injectivity is this monograph's answer to how that question
gets made precise rather than a replacement for it. The observation is
recorded here because a careful reader is likely to notice the
migration independently, and because noticing it is itself a small
piece of evidence that \cref{ch:two-kinds-of-persistence}'s central
move --- replacing \enquote{can this be reconstructed} with
\enquote{does this remain distinguishable} --- was a genuine
sharpening rather than a relabeling: the sharpened primitive is what
the rest of the book actually needed, often enough that the original
naming stopped being where the mathematics lived, even though it
remains where the book's motivation does.

\section{What This Monograph Does Not Establish}

Stating the classification above plainly implies a list of things this
monograph has not shown, and it is worth making the list explicit
rather than leaving it to be inferred.

It has not shown that any real system's designers built retention
architecture with admissibility-relevance in mind, as opposed to the
retention pattern merely correlating with it for unrelated reasons ---
\cref{sec:leak-audit-limits} named this limit directly and it was
never closed. It has not shown that HYDRA's hypothesis-space apparatus
satisfies the descending chain condition \cref{prop:stability}
requires, for HYDRA or for any other real implementation. It has not
shown that any transformer is or is not functionally adequate,
architecturally transparent, both, or neither, over any specific
fiber; \cref{ch:transformer-semantics} built the instrument and did
not use it. It has not surveyed institutional record-keeping, version
control practice, or biological memory research seriously enough to
count as evidence for \cref{ch:general-theory}'s conjecture, only
illustration of it. And it has not resolved the compaction case in
\cref{sec:compaction-ambiguous}, deliberately: that case was left open
because the available evidence does not settle it, not because the
apparatus was insufficient to state the question precisely.

\section{How This Could Be Wrong}

Three specific ways the monograph's non-tautological claims could
fail, stated as concretely as the material allows.

The Admissible Externalization Conjecture fails if a survey of real
systems --- audit trails, version control, institutional record-
keeping, or any other domain satisfying \cref{ch:general-theory}'s
hypotheses --- finds retention patterns statistically uncorrelated
with externally tested admissibility-relevance
(\cref{def:external-relevance}), across enough systems that the
uncorrelated cases cannot be dismissed as measurement noise. This is
exactly \cref{sec:leak-audit-falsification}'s falsification-by-
uncorrelated-retention condition, applied at the scale the conjecture
actually requires rather than to the single system this monograph
audited.

The claim that HYDRA's projection machinery offers a functionally
adequate \emph{and} architecturally transparent account of admissible
computation fails if it turns out, on implementation, that
architectural transparency of the kind \cref{def:architectural-transparency}
describes is not achievable at any scale competitive with the systems
it would need to compete against --- if, in other words, every system
capable of the functional adequacy \cref{ch:transformer-semantics}
describes turns out to purchase it at the cost of exactly the
opacity HYDRA's explicit hypothesis sets were meant to avoid. This
monograph has not ruled this out; \cref{sec:hydra-scope}'s open
questions about fixed versus expanding hypothesis spaces are precisely
where this failure would first appear.

The reconstruction-cost/reconstructive-irreversibility distinction
itself fails as a general-purpose tool, independent of any specific
conjecture built on it, if a significant class of real systems turns
out to have retention patterns that neither classification captures
cleanly --- not because the audit was insufficiently thorough, the way
\cref{sec:compaction-ambiguous}'s cases were merely difficult, but
because a third kind of materialization exists that this monograph's
two-part classification has no room for. Nothing in Part II through
Part IV rules this out; the monograph's strongest claim is only that
the two kinds identified so far are real and distinguishable, not that
they are exhaustive.

\section{Open Problems}

Collected from where they were first raised: whether real hypothesis
spaces satisfy the descending chain condition
(\cref{prop:stability}); whether HYDRA's repair machinery operates
over a fixed or expanding hypothesis space
(\cref{sec:hydra-scope}); how large the gap is between the greedy
allocation heuristic and exact solutions to
\cref{def:allocation-problem}; how the weights and costs that
heuristic requires would be estimated in a real deployment; whether a
genuine replay-level audit of Claude Code, rather than this
monograph's specification-level one, would confirm or overturn any
row of \cref{tab:leak-audit}; whether the compaction case in
\cref{sec:compaction-ambiguous} is admissibility-aligned in the sense
of \cref{def:selective-compression}, which requires access this
monograph did not have; whether \cref{ch:transformer-semantics}'s
protocol, actually run, finds functional adequacy, architectural
transparency, both, or neither, for any specific linguistic fiber; and
whether the literature on constrained Hamiltonian systems has already
developed something resembling this monograph's distinguishability
axis under its own vocabulary, which \cref{ch:variational-primacy}
flagged and did not check; and, now that the check has been made and
found the problem of time and problem of observables to be the closest
existing analogue, whether \cref{rem:problem-of-time-speculation}'s
speculative reframing --- that the dispute persists because a
context-relative question is being asked in a context-free form --- has
any actual purchase on the physics, which would require expertise and
engagement this monograph has not attempted.

\section{Closing}

The original public argument this monograph responds to treated the
discovery of a parser as evidence against understanding. The rebuttal
available from the outset --- that architectural decomposition is not
ontological decomposition, that syntax has long been understood to
possess computational autonomy from semantics --- was sufficient to
dispose of that argument in one chapter, and did not require anything
built afterward. What the remaining chapters pursued was a different
question, one the original debate did not think to ask: not whether a
system decomposes, but whether there is a principled account of what
gets retained when it does, and whether that account can be checked
against a real system rather than merely asserted about one.
Wherever this monograph reached a result it could state as a theorem,
the theorem tended to be close to definitional --- a mark, if the
argument in \cref{ch:falsifiability} was right, not of triviality but
of having found the correct primitive to define. Wherever it reached a
claim that was not definitional --- the leak's actual retention
pattern, the conjecture that the pattern generalizes, the question of
what a transformer's internals actually do --- it said so, and left
the claim exactly as open as the evidence available for it left it. A
theory of what survives reconstruction should not, on its own terms,
claim more permanence for its conclusions than its evidence can
support. This one has tried not to.
% ===== END chapters/ch12-conclusion =====

\appendix
% ===== BEGIN appendices/appA-notation =====
\chapter{Notation and Standing Assumptions}
\label{app:notation}

This appendix collects, in one place, the notation introduced across
Part II and Part III. It is a reference, not an introduction; every
item below is defined properly at its first occurrence in the main
text, cited here by label.

\begin{center}
\renewcommand{\arraystretch}{1.3}
\begin{tabular}{p{2.6cm}p{9.2cm}}
\toprule
\textbf{Symbol} & \textbf{Meaning} \\
\midrule
$X, U, F$ & State space, input space, transition function of a
transition system $(X,U,F)$, $x_{t+1} = F(x_t,u_t)$
(\cref{def:transition-system}). \\
$P$ & A set of admissible predecessor pairs $(x,u) \in X \times U$
under consideration for a given fiber. \\
$\mathrm{Fib}_P(y)$ & The fiber of $y$ under $F$ restricted to $P$:
$\{(x,u) \in P : F(x,u) = y\}$ (\cref{def:fiber}). \\
$s$ & A generating input: a state independent of execution history
from which $X$ is reproducible (\cref{def:generating-input}). \\
$H, G$ & History space and history-update function of a
history-augmented transition system, $h_{t+1} = G(h_t,x_t,u_t)$
(\cref{def:augmented-system}). \\
$\mathrm{GenInput}(X,t)$ & Indicator of whether a generating input for
$X$ remains accessible at time $t$ (\cref{def:time-relative-generating-input}). \\
$\mathcal{X}, C(x)$ & Trajectory space and the continuation family
assigning to each state its permitted continuations
(\cref{def:continuation-family}). \\
$\hat\pi_t, \hat\Pi$ & A system's time-indexed admissibility estimate
and the space of such estimates (\cref{def:projection}); distinguished
from the canonical quotient $\pi^*$. \\
$\pi^*, \mathcal{A}, A^*$ & The canonical admissibility projection, the
hypothesis space of candidate admissibility structures, and the true
structure, per the Minimal Projection Theorem
(\cref{def:hypothesis-space}). \\
$H_t \subseteq \mathcal{A}$ & The hypothesis set surviving after
repair through time $t$ (\cref{def:estimate-as-set}); $H_\infty$ its
eventual stable value (\cref{prop:stability}). \\
$R, \vDash$ & Repair operator and fit relation between estimates and
evidence (\cref{def:repair,def:fit}). \\
$w_i, c_i, \beta$ & Relevance weight, retention cost, and total budget
in the allocation problem (\cref{def:allocation-problem}). \\
$B_t$ & A system's distinction budget at time $t$: collapsed fibers
still recoverable at $t$ (\cref{def:distinction-budget}). \\
$S_G, w_G$ & The subset of fibers over which a compressive transition
$G$ is injective, and its admissibility-alignment status
(\cref{def:selective-compression}). \\
\bottomrule
\end{tabular}
\end{center}

\section{Standing Assumptions}

Two assumptions recur across Part II and Part III and are stated here
together rather than re-derived at each use.

\begin{itemize}
\item \textbf{Well-foundedness where invoked.} Any claim depending on
\cref{prop:stability} (chain stabilization) assumes the relevant
hypothesis space satisfies the descending chain condition. This is
flagged as an assumption, not a theorem, at every point it is used;
\cref{sec:hydra-scope} discusses when it can fail.
\item \textbf{Admissible inputs are exogenously fixed.} The input
space $U$ of a transition system is treated as given rather than
derived; nothing in this monograph offers an account of which inputs
a system \emph{ought} to accept, only of how retention should behave
given a fixed admissibility structure over the inputs it does accept.
\end{itemize}
% ===== END appendices/appA-notation =====
% ===== BEGIN appendices/appB-proofs-part-ii =====
\chapter{Full Proofs from Part II}
\label{app:proofs-part-ii}

The main text proves \cref{thm:historical-distinction} and
\cref{prop:fiber-collapse} in full; this appendix supplies the detail
elided from \cref{prop:non-reducibility}'s proof sketch and from
\cref{rem:chaos-cost}'s chaotic-Hamiltonian illustration, in each case
stating exactly which standard results are being invoked.

\section{Non-Reducibility, in Full}

\begin{proof}[Full proof of \cref{prop:non-reducibility}]
Let $X$ be materialized on reconstruction-cost grounds
(\cref{def:reconstruction-cost}). By \cref{prop:reconstruction-requires-input},
there exists $s$ satisfying \cref{def:generating-input}: $s$ is
independent of the system's execution history and some process $P$
produces $X$ from $s$ without reference to any intermediate state.
Let $\kappa : \mathbb{R}_{\geq 0} \to \mathbb{R}_{\geq 0}$ be any
monotonic cost transformation --- a change in the time, space, or
energy charged for running $P$. Applying $\kappa$ changes the number
$\mathrm{cost}(P, s)$ but does not act on $s$ or on $P$ as
mathematical objects: $s$ remains a state independent of execution
history, and $P$ remains a process taking $s$ to $X$. Hence $s$
continues to satisfy \cref{def:generating-input} after the
transformation, for every $\kappa$, so $X$ remains a reconstruction-
cost case at every cost level. Since \cref{def:irreversibility}
requires the \emph{absence} of a generating input, and $s$'s presence
is invariant under $\kappa$, $X$ never enters the reconstructive-
irreversibility regime under any cost transformation.

Conversely, let a history channel $h$ be materialized on
reconstructive-irreversibility grounds over a fiber
$\mathrm{Fib}_P(y)$ (\cref{def:irreversibility-materialization}), so
by \cref{def:irreversibility} no $s$ satisfying
\cref{def:generating-input} exists for the predecessor's identity.
Suppose, for contradiction, that some cost transformation $\kappa$
brought $X$ into the reconstruction-cost regime. This would require
producing a generating input $s'$ where none existed. But $\kappa$
acts only on the domain of cost values assigned to existing processes;
it is not itself a process, has no output in $X$'s state space, and in
particular is not a state independent of the system's execution
history in the sense \cref{def:generating-input} requires. No
application of $\kappa$ therefore manufactures $s'$. Contradiction; no
such $\kappa$ exists.
\end{proof}

\section{The Chaotic Hamiltonian Case, in Full}

\cref{rem:chaos-cost} asserted that a chaotic Hamiltonian system can
require exponentially growing reconstruction precision while remaining
reconstructively reversible throughout. The precise statement:

\begin{proposition}[Exponential reconstruction cost under positive Lyapunov exponents]
\label{app:prop-lyapunov}
Let $\Phi_t$ be the flow of a Hamiltonian system with largest Lyapunov
exponent $\lambda > 0$. Suppose a reconstruction procedure has access
to $\Phi_t(y)$ only to finite precision $\varepsilon_0$ (measurement or
representation error). Then the precision $\varepsilon(t)$ to which the
predecessor $\Phi_{-t}(y)$ can be recovered by applying $\Phi_{-t}$ to
the imprecise measurement degrades as
\[
\varepsilon(t) \sim \varepsilon_0 \, e^{\lambda t},
\]
so that recovering the predecessor to a fixed target precision
$\varepsilon^*$ requires initial measurement precision
$\varepsilon_0 \sim \varepsilon^* e^{-\lambda t}$, decaying
exponentially in $t$.
\end{proposition}

\begin{proof}[Justification]
This is the standard definition of the largest Lyapunov exponent
applied in reverse time: for nearby trajectories $\Phi_t(y)$ and
$\Phi_t(y) + \delta_0$ with $\delta_0$ small, the Lyapunov exponent
governs $\|\delta_t\| \sim \|\delta_0\| e^{\lambda t}$ under forward
flow for chaotic (positive-exponent) directions. Applying $\Phi_{-t}$
to a measurement perturbed by $\varepsilon_0$ evolves the perturbation
under the time-reversed flow, which exhibits the same exponential
separation in the reversed time direction whenever the forward flow is
chaotic in the corresponding directions, giving
$\varepsilon(t) \sim \varepsilon_0 e^{\lambda t}$ as claimed.
\end{proof}

\begin{corollary}[Cost and irreversibility remain independent under \cref{app:prop-lyapunov}]
\label{app:cor-lyapunov-independence}
\cref{app:prop-lyapunov} exhibits $\varepsilon_0$ (and hence
reconstruction cost, however cost is measured against required
precision) growing without bound as the target precision $\varepsilon^*$
is held fixed and $t \to \infty$, while
\cref{cor:hamiltonian-no-collapse} guarantees $\mathrm{Fib}_P(y)$
remains a singleton for every $t$: the flow map's bijectivity
(\cref{prop:hamiltonian-bijective}) does not depend on $\lambda$ and
holds regardless of whether the system is chaotic. The two
quantities --- required precision and fiber cardinality --- are
governed by unrelated properties of the dynamics (the Lyapunov
spectrum versus the invertibility of $\Phi_t$ as a map), which is the
precise sense in which \cref{prop:non-reducibility} holds for this
example: cost can be driven to infinity without the system ever
entering \cref{def:irreversibility}'s regime.
\end{corollary}
% ===== END appendices/appB-proofs-part-ii =====
% ===== BEGIN appendices/appC-hypothesis-pruning-extended =====
\chapter{The Hypothesis-Pruning Formalism: Extended Convergence Theory}
\label{app:hypothesis-pruning-extended}

\cref{ch:hydra-projections} proved stability under a descending chain
condition (\cref{prop:stability}) without giving a bound on how long
stabilization takes, and proved convergence is not automatic
(\cref{thm:convergence-characterization}) without addressing what
happens when $\mathcal{A}$ is infinite but still well-founded. Both
gaps are closed here.

\section{A Rate Bound for Finite Hypothesis Spaces}

\begin{proposition}[Stabilization bound]
\label{app:prop-stabilization-bound}
If $\mathcal{A}$ is finite, $|\mathcal{A}| = N$, then the hypothesis
sequence $H_0 \supseteq H_1 \supseteq \cdots$ stabilizes after at most
$N - 1$ strict contractions: there exists $T \leq N - 1$ such that
$H_t = H_T$ for all $t \geq T$.
\end{proposition}

\begin{proof}
Each strict contraction $H_t \supsetneq H_{t+1}$ removes at least one
element of $\mathcal{A}$ from consideration, by Contraction
(\cref{def:repair-pruning}). Since $|H_0| \leq N$ and every strict
contraction strictly decreases $|H_t|$ by at least $1$, there can be
at most $N - 1$ strict contractions before $|H_t|$ reaches its
minimum possible value under Soundness, namely $1$ (if the evidence
stream is fully discriminating) or some larger fixed value (if it is
not); in either case no further strict contraction is possible once
$H_t$ stops changing, so stabilization occurs within $N-1$ steps of
the first repair event.
\end{proof}

This bound is tight: a hypothesis space in which each piece of
evidence eliminates exactly one candidate realizes it exactly.

\section{Well-Ordered Infinite Hypothesis Spaces}

\cref{prop:stability} assumed the descending chain condition, which
finite $\mathcal{A}$ satisfies automatically by
\cref{app:prop-stabilization-bound} but infinite $\mathcal{A}$ need
not. A weaker and still useful sufficient condition covers a natural
class of infinite cases.

\begin{proposition}[Stabilization under well-ordering]
\label{app:prop-wellordering}
If there exists a well-ordering $\prec$ of $\mathcal{A}$ such that
Contraction always removes an $\prec$-minimal element of the currently
remaining candidates that is inconsistent with new evidence (a
\enquote{eliminate the least-preferred surviving candidate first}
repair policy), then the hypothesis sequence stabilizes: there is no
infinite strictly decreasing chain, by the defining property of a
well-ordering, that every nonempty subset has a $\prec$-least element,
which rules out an infinite descending chain of distinct nonempty
subsets under this specific elimination policy.
\end{proposition}

\begin{proof}[Justification]
This restates the descending chain condition's content for the
specific case where $\mathcal{A}$ is well-ordered and the repair
policy respects that order: a strictly decreasing chain of subsets of
a well-ordered set, each obtained from the last by removing an
element, corresponds to a strictly decreasing sequence of ordinals
(the order type of what remains), and no strictly decreasing sequence
of ordinals is infinite.
\end{proof}

\begin{remark}
\cref{app:prop-wellordering} is not free: it requires the repair
policy to respect a specific well-ordering, which is an assumption
about how $R$ is implemented, not a fact automatically true of an
arbitrary repair operator satisfying only
\cref{def:repair-pruning}. Repair policies that eliminate candidates
in an order unrelated to any well-ordering of $\mathcal{A}$ are not
covered, and \cref{sec:hydra-scope}'s concern about whether HYDRA's
actual implementation satisfies any such condition is not resolved by
this appendix.
\end{remark}

\section{A Second Non-Convergence Family}

\cref{ex:non-convergence} gave one construction with $|\mathcal{A}| =
2$. The construction generalizes.

\begin{proposition}[Non-convergence persists under enlargement]
\label{app:prop-non-convergence-family}
For any $N \geq 2$, there exists a hypothesis space $\mathcal{A}$ with
$|\mathcal{A}| = N$, a true structure $A^* \in \mathcal{A}$, and an
evidence stream such that $H_t$ stabilizes at $|H_\infty| = N - 1 \geq
1$ without converging, whenever the evidence stream is silent on
exactly one pairwise disagreement among the $N$ candidates.
\end{proposition}

\begin{proof}[Construction]
Let $\mathcal{X} = \{x_1, x_2, x_3\}$ as in \cref{ex:non-convergence},
and let $\mathcal{A} = \{A^*, A_2, \dots, A_N\}$ where $A^*$ equates
$x_1$ and $x_2$, each $A_i$ for $i \geq 3$ disagrees with $A^*$ on some
classification unrelated to $x_1, x_2$ and is eliminated by evidence
addressing that disagreement, and $A_2$ agrees with every $A_i$,
$i\geq 3$ eliminated candidate on all points except $x_1,x_2$, where
it distinguishes them, matching \cref{ex:non-convergence}'s $A'$
exactly. An evidence stream that addresses every disagreement except
the $A^*$-versus-$A_2$ disagreement over $x_1, x_2$ eliminates
$A_3,\dots,A_N$ in finitely many steps (by
\cref{app:prop-stabilization-bound}, since only finitely many
candidates remain to eliminate) and stabilizes at $H_\infty =
\{A^*, A_2\}$, reproducing \cref{ex:non-convergence}'s non-convergence
at this larger scale.
\end{proof}

This shows \cref{ex:non-convergence} is not a boundary artifact of the
smallest possible hypothesis space; non-convergence with stabilization
survives arbitrary enlargement of $\mathcal{A}$, provided the evidence
stream leaves even one pairwise disagreement permanently unaddressed.
% ===== END appendices/appC-hypothesis-pruning-extended =====
% ===== BEGIN appendices/appD-allocation-complexity =====
\chapter{The Allocation Problem: Complexity and Approximation}
\label{app:allocation-complexity}

\cref{ch:general-theory} named \cref{def:allocation-problem} an
instance of weighted knapsack and asserted NP-hardness by citation.
This appendix gives the reduction and proves a concrete approximation
guarantee for the greedy design principle of
\cref{ch:general-theory}, closing the gap left by stating the
heuristic without a bound on how far it can fall short of optimal.

\section{NP-Hardness}

\begin{proposition}[\cref{def:allocation-problem} is NP-hard]
\label{app:prop-nphard}
The decision version of \cref{def:allocation-problem} --- given
$F_1,\dots,F_n$ with weights $w_i$, costs $c_i$, budget $\beta$, and a
target $W$, does a subset $S$ exist with $\sum_{i \in S} w_i \geq W$
and $\sum_{i \in S} c_i \leq \beta$ --- is NP-hard.
\end{proposition}

\begin{proof}[Reduction from 0/1 Knapsack]
The decision version of 0/1 knapsack is: given items with values
$v_i$ and weights $\omega_i$, a capacity $\Omega$, and a target $V$,
does a subset exist with total value at least $V$ and total weight at
most $\Omega$? This is NP-complete (Cormen, Leiserson, Rivest, and
Stein, \emph{Introduction to Algorithms}, standard result). Given any
instance of 0/1 knapsack, construct an instance of
\cref{def:allocation-problem} by setting $n$ equal to the number of
knapsack items, $w_i := v_i$, $c_i := \omega_i$, $\beta := \Omega$,
and $W := V$. A subset $S$ satisfies the allocation problem's
constraints for this instance if and only if it satisfies the
knapsack instance's constraints, by construction; the reduction is
computable in linear time. Since 0/1 knapsack is NP-complete,
\cref{def:allocation-problem} is NP-hard.
\end{proof}

\section{A Proved Approximation Bound}

\cref{ch:general-theory}'s design principle recommended sorting
candidates by $w_i / c_i$ and taking them greedily until the budget is
exhausted. The following makes precise how far this can fall short of
optimal, and gives a simple fix with a proved guarantee.

\begin{definition}[Modified greedy allocation]
\label{app:def-modified-greedy}
Sort candidates so that $w_1/c_1 \geq w_2/c_2 \geq \cdots \geq
w_n/c_n$. Let $k$ be the largest index such that $\sum_{i=1}^{k} c_i
\leq \beta$, and let $S_{\mathrm{greedy}} = \{1,\dots,k\}$. Let $j =
k+1$ be the first candidate that does not fit. The modified greedy
allocation selects $\arg\max\bigl(\textstyle\sum_{i \in
S_{\mathrm{greedy}}} w_i,\ w_j\bigr)$: whichever of the greedy prefix or
the single next-best candidate has higher total weight.
\end{definition}

\begin{theorem}[Modified greedy is a 2-approximation]
\label{app:thm-2-approx}
Let $\mathrm{OPT}$ denote the optimal value of
\cref{def:allocation-problem} for a given instance, and let
$\mathrm{ALG}$ denote the value achieved by
\cref{app:def-modified-greedy}. Then $\mathrm{ALG} \geq
\mathrm{OPT}/2$.
\end{theorem}

\begin{proof}
Consider the fractional relaxation of the problem, in which each
candidate may be partially included. The optimal fractional solution
$\mathrm{OPT}_f$ is achieved by taking candidates in decreasing order
of $w_i/c_i$ --- exactly the sorted order above --- filling the budget
completely, taking a fractional amount of the first candidate that
does not fit whole. This is a standard property of the fractional
knapsack problem: greedy-by-ratio is optimal for the fractional
relaxation. The fractional optimum therefore equals $\sum_{i \in
S_{\mathrm{greedy}}} w_i + \theta w_j$ for some $\theta \in [0,1)$,
where $S_{\mathrm{greedy}}$ and $j$ are as in
\cref{app:def-modified-greedy}. Since $\mathrm{OPT} \leq
\mathrm{OPT}_f$ (the integer optimum cannot exceed the fractional
relaxation's optimum), and $\mathrm{OPT}_f \leq \sum_{i \in
S_{\mathrm{greedy}}} w_i + w_j$ (bounding $\theta < 1$ by $1$), we have
\[
\mathrm{OPT} \leq \sum_{i \in S_{\mathrm{greedy}}} w_i + w_j \leq 2 \max\Bigl(\sum_{i \in S_{\mathrm{greedy}}} w_i,\ w_j\Bigr) = 2 \cdot \mathrm{ALG}.
\]
Dividing by $2$ gives $\mathrm{ALG} \geq \mathrm{OPT}/2$.
\end{proof}

\begin{remark}[What this adds to \cref{ch:general-theory}]
\cref{app:thm-2-approx} upgrades the design principle from a heuristic
with no stated guarantee to one with a proved worst-case bound: a
designer following \cref{app:def-modified-greedy} --- the ordinary
greedy allocation, augmented with a single check against the best
individual candidate --- never does worse than half of what an
omniscient allocator could achieve, regardless of how the weights and
costs are distributed. This does not close
\cref{sec:general-open}'s first gap entirely, since a factor of two is
often too coarse a guarantee for a real deployment decision, but it
replaces \enquote{a heuristic with an unknown gap} with \enquote{a
heuristic with a proved gap of at most this size,} which is the
precise sense in which this appendix answers the question
\cref{sec:general-open} left open rather than merely restating it.
\end{remark}
% ===== END appendices/appD-allocation-complexity =====
% ===== BEGIN appendices/appE-hamiltonian-flow =====
\chapter{Hamiltonian Flow: Existence, Uniqueness, and Liouville's Theorem}
\label{app:hamiltonian-flow}

\cref{ch:hamiltonian-lagrangian-continuation} cited standard ODE
theory for \cref{prop:hamiltonian-bijective} and stated Liouville's
theorem without derivation. Both are supplied here.

\section{Existence, Uniqueness, and the Flow Map}

\begin{theorem}[Picard--Lindel\"of / Cauchy--Lipschitz]
\label{app:thm-picard-lindelof}
Let $f : \mathbb{R}^n \times \mathbb{R} \to \mathbb{R}^n$ be Lipschitz
continuous in its first argument, uniformly in the second, on a
neighborhood of $(z_0, t_0)$. Then the initial value problem $\dot z =
f(z,t)$, $z(t_0) = z_0$ has a unique solution on some interval
containing $t_0$.
\end{theorem}

This is classical and not reproved here; it is the standard existence-
uniqueness theorem for ordinary differential equations. For a
Hamiltonian system, $z = (q,p)$ and $f(z,t) = (\partial H/\partial p,
-\partial H/\partial q)$, Lipschitz whenever $H$ is $C^2$.

\begin{proof}[Proof of \cref{prop:hamiltonian-bijective}]
Fix $t$. Uniqueness (\cref{app:thm-picard-lindelof}) gives that
distinct initial conditions $z_0 \neq z_0'$ produce distinct solutions
$\Phi_t(z_0) \neq \Phi_t(z_0')$ for as long as both solutions are
defined: if $\Phi_t(z_0) = \Phi_t(z_0')$ for some $t$, then running the
unique solution through $(\Phi_t(z_0), t)$ backward to time $t_0$
would have to recover both $z_0$ and $z_0'$, contradicting uniqueness
unless $z_0 = z_0'$. This gives injectivity of $\Phi_t$. Surjectivity
onto its domain of definition follows by running the flow forward from
every point; smoothness of $\Phi_t$ and of its inverse $\Phi_{-t}$
follows from smooth dependence of solutions of an ODE on initial
conditions, a standard strengthening of
\cref{app:thm-picard-lindelof} under $C^1$ (here $C^2$, since $H \in
C^2$) hypotheses on $f$.
\end{proof}

\section{Liouville's Theorem}

\begin{theorem}[Liouville]
\label{app:thm-liouville}
Let $\Phi_t$ be the flow of a Hamiltonian vector field $X_H =
(\partial H/\partial p, -\partial H/\partial q)$ on phase space $M$.
Then $\Phi_t$ preserves the standard volume form: for any measurable
region $\Omega \subseteq M$, $\mathrm{vol}(\Phi_t(\Omega)) =
\mathrm{vol}(\Omega)$ for all $t$.
\end{theorem}

\begin{proof}
By the divergence theorem applied to the flow of a vector field, the
rate of change of volume under the flow of $X_H$ is governed by the
divergence of $X_H$:
\[
\frac{d}{dt}\mathrm{vol}(\Phi_t(\Omega)) = \int_{\Phi_t(\Omega)} \nabla \cdot X_H \, d\mathrm{vol}.
\]
Compute the divergence directly:
\[
\nabla \cdot X_H = \sum_i \left( \frac{\partial}{\partial q_i}\frac{\partial H}{\partial p_i} + \frac{\partial}{\partial p_i}\left(-\frac{\partial H}{\partial q_i}\right) \right) = \sum_i \left( \frac{\partial^2 H}{\partial q_i \partial p_i} - \frac{\partial^2 H}{\partial p_i \partial q_i} \right) = 0,
\]
by equality of mixed partial derivatives (Clairaut's theorem, given
$H \in C^2$). Since the divergence vanishes identically, the rate of
change of volume is zero for every $\Omega$ and every $t$, so volume
is preserved.
\end{proof}

\begin{remark}[Relation to \cref{cor:hamiltonian-no-collapse}]
\cref{app:thm-liouville} is the continuous, measure-theoretic
strengthening of \cref{cor:hamiltonian-no-collapse}'s pointwise claim.
Bijectivity of $\Phi_t$ (\cref{prop:hamiltonian-bijective}) says no
two points are ever identified by the flow; Liouville's theorem says,
in addition, that the flow does not even locally compress or
rarefy the space of possibilities --- volume in phase space is neither
created nor destroyed. Neither statement requires the system to be
integrable or non-chaotic; both hold identically for the chaotic
system considered in \cref{app:prop-lyapunov}, which is precisely why
\cref{app:cor-lyapunov-independence} could assert unbounded
reconstruction cost coexisting with zero information loss --- volume
preservation guarantees the information is never destroyed, while
saying nothing about how finely that preserved volume becomes folded
and stretched, which is exactly what governs reconstruction cost.
\end{remark}
% ===== END appendices/appE-hamiltonian-flow =====
% ===== BEGIN appendices/appF-dirac-bergmann =====
\chapter{The Dirac--Bergmann Algorithm}
\label{app:dirac-bergmann}

\cref{ch:variational-primacy} used the vocabulary of primary
constraints, gauge orbits, and Dirac observables without deriving the
algorithm that produces them. This appendix gives the construction in
outline, precisely enough to support the claims
\cref{rem:gauge-vs-admissibility} makes about its logical form.

\section{Primary Constraints from a Singular Hessian}

Let $L(q,\dot q)$ be a Lagrangian on configuration space with
coordinates $q^i$, $i = 1,\dots,n$. The Legendre transform defines
momenta $p_i := \partial L/\partial \dot q^i$. The transform is
invertible for $\dot q^i$ in terms of $(q,p)$ exactly when the Hessian
\[
W_{ij} := \frac{\partial^2 L}{\partial \dot q^i \partial \dot q^j}
\]
is non-singular. When $\det W = 0$, the Hessian has a kernel of some
dimension $m > 0$: there exist $m$ independent null vectors
$v^i_{(a)}$, $a = 1,\dots,m$, with $W_{ij} v^j_{(a)} = 0$. Each null
vector generates a relation among the momenta not involving the
velocities, a \emph{primary constraint}
\[
\phi_a(q,p) \approx 0, \qquad a = 1,\dots,m,
\]
where $\approx$ (weak equality) denotes equality only on the
constraint surface, not as an identity on the full phase space --- the
distinction Dirac's formalism depends on throughout, since quantities
that vanish weakly may still have nonzero Poisson brackets with other
phase space functions.

\section{Consistency and Secondary Constraints}

The canonical Hamiltonian $H_c = p_i \dot q^i - L$ is well-defined on
the constraint surface despite the non-invertibility, since the
velocity-dependence of $p_i\dot q^i - L$ drops out in directions along
$\ker W$. The total Hamiltonian is $H_T = H_c + \lambda^a \phi_a$,
with $\lambda^a$ undetermined Lagrange multipliers. Consistency of the
primary constraints under time evolution, $\dot\phi_a = \{\phi_a, H_T\}
\approx 0$, either determines some of the $\lambda^a$, is
automatically satisfied, or generates a new, \emph{secondary}
constraint $\phi_a' \approx 0$ not implied by the primary constraints
alone. The procedure iterates --- checking consistency of each new
constraint --- until no further constraints are generated; this is the
Dirac--Bergmann algorithm.

\section{First- and Second-Class Constraints}

Constraints are classified by the Poisson bracket algebra they
generate on the constraint surface.

\begin{definition}[First- and second-class constraints]
A constraint $\phi$ (primary or secondary) is first-class if its
Poisson bracket with every other constraint vanishes weakly:
$\{\phi, \phi_b\} \approx 0$ for all constraints $\phi_b$. A
constraint that is not first-class is second-class.
\end{definition}

First-class constraints are, in the standard interpretation this
monograph's \cref{ch:variational-primacy} follows, the generators of
gauge transformations: the transformation $\delta z = \epsilon
\{z,\phi\}$ generated by a first-class constraint $\phi$ maps points
on the constraint surface to other points on the constraint surface
satisfying the same equations of motion, and is standardly interpreted
as relating physically identical states --- the gauge orbit structure
\cref{rem:gauge-vs-admissibility} discusses.

Second-class constraints, by contrast, do not generate gauge
transformations; they are eliminated by passing to the Dirac bracket,
\[
\{F,G\}^* := \{F,G\} - \{F,\chi_m\}(M^{-1})^{mn}\{\chi_n,G\},
\]
where $\chi_m$ ranges over the second-class constraints and $M^{mn} :=
\{\chi_m,\chi_n\}$ is invertible on the second-class constraint
surface by definition of second-class. The Dirac bracket restricts
dynamics to the surface where the second-class constraints hold
identically.

\section{Dirac Observables}

\begin{definition}[Dirac observable]
A phase space function $\mathcal{O}$ is a Dirac observable if
$\{\mathcal{O}, \phi_a\} \approx 0$ for every first-class constraint
$\phi_a$: it is invariant under every gauge transformation the theory
admits.
\end{definition}

This is the criterion \cref{rem:gauge-vs-admissibility} contrasted
directly with \cref{def:admissibility-relevant}: membership is decided
by an algebraic condition against the full set of first-class
constraints, fixed once the theory is fixed, with no reference to any
downstream task, continuation family, or future admissibility
judgment. \cref{ch:variational-primacy}'s claim that this differs in
kind from \cref{def:admissibility-relevant}'s relativized criterion
rests on exactly this contrast: $\{\mathcal{O},\phi_a\}\approx 0$ is a
single computation with a single answer, while
\cref{def:external-relevance}'s counterfactual test is parameterized
by a choice of continuation family that need not be, and in the
disputed general-relativistic case discussed in
\cref{ch:variational-primacy} demonstrably has not been, agreed upon.

\section{The Free Particle Example, Condensed}

The reparametrization-invariant free relativistic particle, used in
\cref{ch:variational-primacy}'s source material as a worked
illustration, has Lagrangian $L = \tfrac{1}{2n}\dot q^2 - \tfrac{1}{2}n$
with $q^\mu$ the particle's spacetime coordinates and $n$ an auxiliary
lapse variable. The momentum conjugate to $n$ vanishes identically,
$\pi = \partial L/\partial \dot n = 0$, giving one primary constraint
directly. Consistency of $\pi \approx 0$ under time evolution generates
the secondary constraint $p^2 + 1 \approx 0$ (the mass-shell
condition), and the algorithm terminates: both constraints are
first-class, since their mutual Poisson bracket vanishes, and together
they generate the reparametrization gauge freedom of the model. This
is the smallest nontrivial instance of the structure
\cref{ch:variational-primacy} scales up to general relativity, where
the analogous constraints are the Hamiltonian and momentum constraints
whose interpretation --- pure gauge, or partly physical --- is the
substance of the dispute \cref{rem:gauge-vs-admissibility} and
\cref{rem:problem-of-time-speculation} describe.
% ===== END appendices/appF-dirac-bergmann =====

\backmatter
\cleardoublepage
\addcontentsline{toc}{chapter}{Bibliography}
% ===== BEGIN bibliography =====
\begin{thebibliography}{99}

\bibitem{anderson2017}
Anderson, E. (2017). \emph{The Problem of Time: Quantum Mechanics
Versus General Relativity}. Fundamental Theories of Physics, vol.\
190. Cham: Springer.

\bibitem{arnold1989}
Arnold, V.\ I. (1989). \emph{Mathematical Methods of Classical
Mechanics} (2nd ed.). New York: Springer.

\bibitem{bergmann-komar1960}
Bergmann, P.\ G., \& Komar, A. (1960). Poisson brackets between
locally defined observables in general relativity. \emph{Physical
Review Letters}, 4, 432--433.

\bibitem{chomsky1957}
Chomsky, N. (1957). \emph{Syntactic Structures}. The Hague: Mouton.

\bibitem{chomsky1965}
Chomsky, N. (1965). \emph{Aspects of the Theory of Syntax}.
Cambridge, MA: MIT Press.

\bibitem{cormen2009}
Cormen, T.\ H., Leiserson, C.\ E., Rivest, R.\ L., \& Stein, C.
(2009). \emph{Introduction to Algorithms} (3rd ed.). Cambridge, MA:
MIT Press.

\bibitem{cover-thomas2006}
Cover, T.\ M., \& Thomas, J.\ A. (2006). \emph{Elements of
Information Theory} (2nd ed.). Hoboken, NJ: Wiley.

\bibitem{dirac1964}
Dirac, P.\ A.\ M. (1964). \emph{Lectures on Quantum Mechanics}.
Belfer Graduate School of Science Monographs Series No.\ 2. New York:
Yeshiva University.

\bibitem{goswami2026}
Goswami, M. (2026, July 6). \emph{Claude Code, Chomsky, and the
Illusion of the Wizard Behind the Curtain}. LinkedIn.
\url{https://www.linkedin.com/pulse/claude-code-chomsky-illusion-wizard-behind-curtain-mirtunjaya-goswami-qhr5c/}

\bibitem{kuchar1992}
Kucha\v{r}, K.\ V. (1992). Time and interpretations of quantum
gravity. In G.\ Kunstatter, D.\ E.\ Vincent, \& J.\ G.\ Williams
(Eds.), \emph{Proceedings of the 4th Canadian Conference on General
Relativity and Relativistic Astrophysics}. Singapore: World
Scientific. Reprinted in \emph{International Journal of Modern
Physics D}, 20 (Supp.\ 1), 3--86 (2011).

\bibitem{lextrait2026}
Lextrait, V. (2026). \emph{Behind Claude Code's Curtain: Generative
Grammars Doing Heavy Lifting}. LinkedIn. Cited as reported in
Goswami (2026); the original post's exact wording could not be
independently verified for this monograph.

\bibitem{mohallem2024}
Mohallem, J.\ R. (2024). \emph{Lagrangian and Hamiltonian Mechanics:
A Modern Approach with Core Principles and Underlying Topics}. Cham:
Springer.

\bibitem{pons-salisbury2005}
Pons, J.\ M., \& Salisbury, D.\ C. (2005). The issue of time in
generally covariant theories and the Komar--Bergmann approach to
observables in general relativity. arXiv:gr-qc/0503013.

\end{thebibliography}

\end{document}

